% ==================================================================
% ch01 -- from the research note "The Mortal Scientist"
% ==================================================================
\chapter{The Mortal Scientist}
\label{ch:sci}

\noindent
The preceding parts asked what a physics is and why a scientist inside one
would be misled. Both took the scientist for granted. This chapter builds one.

The constraint that does the work is stated in a single sentence: a
computation can only find out about another computation by interacting with
it, and interacting costs. Everything else follows from taking that
seriously and refusing to add anything not already forced. There is no
dataset, because a dataset is somebody else's observations supplied for
free. There is no scorekeeper, because a scorekeeper would be an authority
standing outside the world telling the learner how it is doing, and no such
authority is available across the cut. What is left is a term in parallel
with other terms, holding a finite supply of tokens, spending them to look,
and dying when it runs out.

Three things then have to be supplied and only one of them turns out to be a
choice. The environment is fixed by ontological isolation. The hypothesis
language is not chosen but \emph{generated} --- fed the presentation of the
calculus, the OSLF construction emits a logic whose formulae denote exactly
the equivalence classes the learner could distinguish by experiment, which is
the Pledge's ``anything Bob can do Alice can do'' turned into syntax. The
motive is the one genuinely free decision, and we make it metabolic: the
learner that predicts badly cannot afford to look, and the learner that
cannot look starves.

What comes out is stranger than what went in. Destructive measurement forces
induction. Foraging bounds concealment, so that perfect privacy is available
only to something already dying. Reproduction turns out to be reflection ---
a genome is a quotation of code, inert because quotation does not reduce ---
and sex is a defence against a predator's hypothesis language rather than a
mechanism of variation. And the ecology divides, without being told to, into
the two kinds of organism Eric Smith finds in the biosphere: those that eat
only external energy, and those that eat each other.

\bigskip

\section{Introduction}\label{sci:S1}

\subsection{The scientist as a learner}\label{sci:S1.1}

Most computational accounts of learning begin with a dataset: a supply of observations that exists
independently of the learner, whose acquisition is free, and whose acquisition does not disturb the
source. Given such a supply the problem is one of induction --- find the hypothesis that best accounts
for what is already in hand --- and the interesting questions are statistical.

A scientist has none of this. There is no dataset waiting; there is a world, and the scientist must
decide what to do to it in order to make it answer. The characteristic activity of the scientific
method is therefore not induction but \emph{the design of experiments}: choosing an intervention whose
outcome would discriminate between the hypotheses currently in play, performing it, and letting the
result decide. Induction is what happens afterwards, and it is the easier half.

Three features of that activity are ordinarily abstracted away and are exactly what we want to keep.

\begin{description}
\item[Observation is action.] To learn something about a system one must do something to it. There is
no passive channel. Every measurement is an intervention, and the system afterwards is not the system
before.
\item[Inquiry is expensive.] Experiments consume resources that could have been spent otherwise, and a
scientist with a finite endowment can run only finitely many. Which questions are worth asking is
therefore not a philosophical question but a budgetary one.
\item[The scientist is in the world it studies.] It is not a disembodied observer looking on from
outside. It is made of the same stuff as its subject matter, subject to the same constraints, and
visible to whatever else is looking.
\end{description}

A model that keeps all three will look less like statistics and more like ecology, and that is the
model we build.

\subsection{What such a model must supply}\label{sci:S1.2}
\label{sci:sec:requirements}

Before any formalism, here is what has to be on the table.

\begin{description}
\item[\textnormal{(R1)} An environment, and a ceiling on what can be known of it.] One must say what
the learner is learning \emph{about}, and --- crucially --- what the best possible outcome of inquiry
would be. Without a ceiling there is no way to say whether a learner is doing well.
\item[\textnormal{(R2)} A language for hypotheses.] Hypotheses must be stated in something. That
something should be able to express every distinction the learner could in principle observe, and no
more. A language too coarse leaves the learner unable to say what it can see; a language too fine
sends it chasing differences that no experiment could ever settle.
\item[\textnormal{(R3)} A mechanism from hypothesis to verdict.] There must be an account of what an
experiment \emph{is}, how a hypothesis determines one, and how the outcome comes back as confirmation
or refutation --- including an honest account of the case where nothing comes back at all.
\item[\textnormal{(R4)} A motive.] Something must make the learner prefer better hypotheses. If this
is an external scorekeeper handing out rewards, the model has outsourced the very thing it set out to
explain, and one may reasonably ask who scores the scorekeeper.
\end{description}

\subsection{Four answers}\label{sci:S1.3}

Our answers are as follows, and the remainder of the introduction says why each is less a choice than
it appears.

\paragraph{(R1) The environment is other computations.} We take as premise that computations are
ontologically isolated: a computation has no access to anything but through interaction. It follows
immediately that the environment of a computation is other computations, and that what one can come to
know of another is bounded above by the equivalence induced by the contexts one can build against it
--- context-labelled bisimulation \cite{Probing}. That is the ceiling R1 demands, and it is derived
rather than posited. We fix the model of computation as the \rhoc\ calculus \cite{RHO}.

\paragraph{(R2) The hypothesis language is generated, not designed.} Given the ceiling, R2 has a
canonical answer: the language must be adequate for context-labelled bisimulation, which is exactly
what a logic generated from the calculus' own presentation gives. Feeding the \rhoc\ calculus'
presentation to the OSLF family of algorithms \cite{OSLF} emits namespace logic
\cite{NamespaceLogic}, whose formulae denote bisimulation equivalence classes. We admit the whole of
it, including the structural predicates --- and admitting those is what gives the learner senses
(\S\ref{sci:sec:logic}).

\paragraph{(R3) An assay is a guarded receipt.} We work in the variant of the \rhoc\ calculus whose
for-comprehension carries a \texttt{where} clause \cite[\S9]{Omnibus}, so a formula may sit in guard
position on a receipt. A hypothesis installed as a guard \emph{is} an experiment, and a confirmation
is a rendezvous. The checker is the reduction rule; nothing is interpreted, and the test is metered by
the same apparatus that meters everything else.

\paragraph{(R4) The motive is metabolic.} Cost accounting as a monad \cite{CostMonad} makes the
learner mortal: every step is paid for, and a computation that cannot pay stops. We then gamify.
Tokens are conserved and never minted; the learner holds a bounded internal supply; further stacks lie
distributed through the environment, including inside other computations; and breaking a computation
open to harvest its stack is the analogue of eating. This is the answer to the scorekeeper problem. A
hypothesis says \emph{where the tokens are}, and the reward for a good one is the tokens themselves.
Predictive accuracy and metabolic success become the same quantity, measured in the same units, with
nobody adjudicating.

\medskip
\noindent That is the whole of what we put in. The rest of the note is largely the discovery that we
had no choices left: \S\ref{sci:sec:audit} audits which features are forced by this list and which remain
free, and the list of free parameters is short.

\subsection{What we are claiming, and what we are not}\label{sci:S1.4}

We are claiming that the commitments just listed determine almost all of the apparatus one would
otherwise have to invent. The interesting sections below are the ones in which we find we had no
choices left.

We do treat learning, and \S\ref{sci:sec:search} is where. What we give there is the geometry of the
hypothesis space and the constraints any revision policy must satisfy in it --- that the space is an
ultrametric tree and so admits no gradient, that incremental revision provably cannot leave the ball it
begins in, that the cost of a move is monotone in its distance, and that the branching a searcher can
survey is bounded by its wealth --- together with a catalogue of strategies compatible with those
constraints. We are \emph{not} claiming to have fixed a policy, proved a convergence theorem, or said
anything about statistical efficiency. Which policy is best is an ecological question and depends on
the distribution of \S\ref{sci:sec:distributions}; the policy itself remains one of only three genuinely
free parameters we find. We do treat reproduction (\S\ref{sci:sec:repro}), but we do not treat inheritance of
acquired characteristics: whether a lineage's transmitted material can carry what an individual
learned is left open, since the empirical case against a strict germ--soma barrier
\cite{Noble} is strong enough that we would rather not build on either answer.

We are also not claiming a novel logic. The logic is the one obtained for free from the encoding, in
exactly the sense of \cite{ClassifierNote}: given the presentation $(\Sigma, E, R)$ of the \rhoc\
calculus, the OSLF algorithms emit a logic, and it is the same logic whether one is classifying knot
diagrams or hunting for tokens. What is new here is the reading, the game, and three consequences of
the game that we did not anticipate.

\subsection{Why adequacy is the right demand}\label{sci:S1.5}

R2 asked for a language neither too coarse nor too fine, and it is worth saying why that pair of
failures is the right pair to worry about, since adequacy is a technical condition that can look like
a formality.

A learner separated from its environment by an interaction cut has exactly one channel of access: what
interaction reveals. If its hypothesis language is too coarse, there are distinctions it can observe
but cannot state, and it is blind in a way its own faculties do not require. If the language is too
fine, it will form hypotheses that differ on environments no experiment could separate, and will spend
tokens --- which by R4 are its life --- discriminating between what are, for it, the same world.
Adequacy rules out both at once: the formulae separate precisely what interaction separates.

This is why the logic must be generated from the same presentation that generates the calculus rather
than designed. A designed logic has no reason to be adequate. A generated one is adequate whenever its
target specification carries the connectives adequacy needs \cite[\S1.3]{ClassifierNote}, and when it
does not, the failure is legible and located.

\begin{remark}[Adequate for which relation]
\label{sci:rem:adequacy-index}
Said that baldly, the demand is ambiguous, and the ambiguity is worth removing before it does damage
later in this chapter. ``What interaction separates'' is not a single relation. The generated logic has
a structural layer as well as a modal one, and a decomposition is not a step, so the structural layer
sees things context-labelled bisimulation in the base theory does not. Chapter~\ref{ch:obsext} settles
this: observational equivalence is indexed by the destructuring instruments admitted to the observer,
and the relation for which the generated logic is adequate is bisimulation in the observer extension
$\Obs(\GSLT,\Ob)$ for a stated observation signature $\Ob$
(Corollary~\ref{ox:cor:adequacy}). Everywhere this chapter says the formulae separate exactly what
interaction separates, read: exactly what interaction separates \emph{for a learner holding these
instruments}. That is a weaker claim than it looks, and it is the one that is true.
\end{remark}

\subsection{Ideal logic, checkable fragment, mortal scientist}\label{sci:S1.6}

We inherit the methodological commitment of \cite{ClassifierNote}: the \emph{ideal} logic, adequate
for the idealised calculus, is what defines the invariants; \emph{fragments}, restricted so that model
checking terminates, are what decide them. In the present setting this distinction stops being
methodology and becomes the subject matter. A mortal scientist cannot afford the ideal logic. What
fragment it can afford is constrained by its budget, and that constraint is a large part of its
inductive bias. We return to this in \S\ref{sci:sec:lattice}, and it is the reason this chapter's
central objects are approximants rather than limits.

Two levels are one too few, and the missing one is not the budget.

\begin{remark}[Three restrictions, and only the third is a budget]
\label{sci:rem:three-limits}
Following Remark~\ref{ox:rem:threerestrictions}, a learner's reach is limited in three separable ways.
The \emph{ideal observer} holds the full observation signature and reaches $=_{\eqs}$: this is the
level at which adequacy is asserted. The \emph{capability-limited} observer holds only some
$\Ob \subsetneq \Sigma$ of the destructuring instruments --- it can take some things apart and not
others --- and its reach is the corresponding point of the family of
Proposition~\ref{ox:prop:mono}. The \emph{budget-limited} observer holds instruments it cannot afford
to use, and its reach is smaller again.

Which instruments a scientist holds and what it can afford to run are different facts about it, and
they answer to different pressures: the first is a fact about the world it was born into and the
extension its presentation admits, the second about its stack. This chapter is mostly about the third
restriction, which is the novel one. It is not entitled to the conclusion that the third is the
\emph{only} one, and where the argument below reads as though budget alone fixes the accessible
hypothesis class, the second restriction has been silently held constant.
\end{remark}

\section{Isolation, the cut, and what a computation can know}\label{sci:S2}
\label{sci:sec:cut}

\subsection{The premise}\label{sci:S2.1}

\begin{quote}
Computations are ontologically isolated. Therefore the environment of a computation is other
computations.
\end{quote}

We take this as a premise rather than argue for it. Its immediate formal content is that a world
factors across an interaction cut. In an interactive graph-structured lambda theory \cite{Omnibus},
an interaction cut $K$ separates a program from its environment: application for the lambda calculus,
parallel composition for CCS, the pi calculus, and the \rhoc\ calculus. We write
\[
W \;\equiv\; S \mid E
\]
for a world factored into a scientist $S$ and its environment $E$. Nothing here privileges $S$: the
factorisation is a choice of which side one is standing on, and in \S\ref{sci:sec:mixed} we will take
seriously that $E$ contains further scientists for whom $S$ is environment.

\subsection{Three grades of access, and only three}\label{sci:S2.2}

The premise has a sharp consequence that the rest of the note leans on. $S$ does not hold $E$'s term.
It cannot read $E$'s structure by inspection, because inspection is not an operation the calculus
provides across a cut. What it has instead is exactly three things, and it is worth being explicit
that the calculus offers no fourth.

\begin{description}
\item[Perception.] $S$ may evaluate structural predicates at the \emph{interaction surface} --- the
outermost interaction structure of $E$, the part exposed at the cut. This is what senses are: one
sees the skin of a rock, not its interior.
\item[Interaction.] $S$ may build a probe $C[-]$, plug $E$ into it, and observe the steps that
$C[E]$ can then take. This is the only route to behaviour. It is worth being careful here, because
two different contexts are in play and the note has more than one use for each. The probe $C$ is
what $S$ constructs and is a context \emph{around} $E$. The label on the generated modality
$\langle K_j \rangle$ of \S\ref{sci:sec:logic} is a context around the \emph{redex} --- everything in
$C[E]$ except the interacting pair --- so it includes $S$'s own remaining code and whatever of $E$
is not participating. The two coincide exactly when the redex sits at the cut, which is the case
$S$ can engineer and the reason probing is informative at all; they come apart as soon as $S$
attends to a step internal to $E$. We write $\langle K_j \rangle$ throughout for the generated
modality and reserve $C[-]$ for probes.
\item[Reflection.] If $S$ \emph{holds} a name, it may look inside it, because in a reflective calculus
a name is a quoted process and the name predicate $@\varphi$ sees into that quotation without
communicating at all.
\end{description}

The third is the interesting one and it is available only for names $S$ has come to hold. This gives
the distinction the note turns on: \textbf{observation versus disclosure}. Behaviour is what one
observes across a cut; structure is what one is given, or takes. Perception reads the surface of
things one does not hold; reflection reads the interior of things one does. The operation converting
the second into the first is the subject of \S\ref{sci:sec:game}, and it is fatal to the thing converted.

\begin{figure}[t]
\centering
\fitwidth{%
\begin{tikzpicture}[>=Stealth, font=\small]
  % scientist
  \draw[rounded corners=3pt, thick] (-4.6,-1.5) rectangle (-1.4,1.5);
  \node at (-3.0,1.15) {\bfseries $S$ (scientist)};
  \node at (-3.0,0.45) {internal stack $\sigma_S$};
  \node at (-3.0,-0.15) {hypothesis $\varphi$};
  \node at (-3.0,-0.75) {revision policy};
  % environment
  \draw[rounded corners=3pt, thick] (1.4,-1.9) rectangle (5.4,1.9);
  \node at (3.4,1.55) {\bfseries $E$ (environment)};
  \node at (3.4,0.85) {other scientists};
  \node at (3.4,0.25) {mortal computations};
  \node at (3.4,-0.35) {bare token stacks};
  \node at (3.4,-1.25) {\itshape interior: opaque};
  % the cut
  \draw[very thick, dashed] (-1.0,-2.3) -- (-1.0,2.3);
  \node[rotate=90, anchor=south] at (-1.15,0) {\scriptsize interaction cut $K \;=\; \mid$};
  % surface
  \draw[thick] (1.4,-1.9) -- (1.4,1.9);
  \node[rotate=90, anchor=north] at (1.55,0) {\scriptsize surface};
  % arrows
  \draw[->, thick] (-1.4,1.05) to[out=20,in=160] node[above,pos=.5]{\scriptsize perceive} (1.4,1.05);
  \draw[->, thick] (-1.4,0.05) to[out=15,in=165] node[above,pos=.5]{\scriptsize assay} (1.4,-0.35);
  \draw[->, very thick] (-1.4,-1.05) to[out=-15,in=195] node[below,pos=.5]{\scriptsize break in} (1.4,-1.25);
\end{tikzpicture}}
\caption{The cut. Perception evaluates structural predicates at the surface; an assay reaches across
and perturbs; a break-in passes the surface and is fatal to what lies behind it. There is no fourth
route, and the three are ordered by both what they reveal and what they destroy
(\S\ref{sci:sec:exposure}).}
\label{sci:fig:cut}
\end{figure}

\section{The logic, and senses}\label{sci:S3}
\label{sci:sec:logic}

\subsection{What the generator supplies}\label{sci:S3.1}

We recall only what we use, following \cite[\S2]{ClassifierNote}. The \rhoc\ calculus has term formers
$0$, $\ast n$, $n!(Q)$, $\texttt{for}(x \leftarrow n)P$, parallel composition, and $@P$; parallel
composition carries an associative--commutative equational theory and there is a single communication
rewrite. Feeding this presentation to the generator yields:

\paragraph{Structural connectives.} One connective per term former. We write $\mathrm{out}(n)\varphi$
for the predicate holding of $n!(Q)$ with $Q \models \varphi$, and $\mathrm{in}(n)\varphi$ for the
analogue on receipts. Because $\mid$ is associative--commutative, the generated connective for
parallel composition is a separating conjunction:
\[
P \models \varphi \mid \psi \iff P \equiv Q \mid R \text{ with } Q \models \varphi,\; R \models \psi .
\]

\paragraph{Behavioural modalities.} One primitive modality $\langle K_j \rangle$ per rewrite rule and
choice of redex position, labelled by the one-hole context in which the step fires. The index is not
decoration and we keep it. The rho calculus has a single interaction rule, so all the information in
a label is in the \emph{position}: for a world $W \equiv K_j[\,\texttt{for}(x \leftarrow n)P \mid
n!(Q)\,]$ the label $K_j$ is everything else. Consequently $\exists j.\,\langle K_j\rangle\varphi$
and $\forall j.\,[K_j]\varphi$ quantify over the redex positions of $W$, which is what gives the
modal layer its force: an action-labelled logic sees one label where this one sees as many as there
are enabled interactions. Where a formula below is written with a bare $\langle K\rangle$ it
abbreviates the existential, and where the quantifier is doing work we write it out.

\paragraph{Name predicates.} A name is a quoted process, so
\[
n \models @\varphi \iff n = @P \text{ with } P \models \varphi ,
\]
which is namespace logic \cite{NamespaceLogic}: one describes a namespace by a property rather than
enumerating it. \S\ref{sci:sec:forced} shows this is not a convenience but a requirement.

\paragraph{Restriction, fixed point, grading.} The hidden-name quantifier
$P \models \mathsf{H}x.\varphi$ iff $P \equiv \texttt{new } x \texttt{ in } Q$ with
$Q \models \varphi$ and $x$ fresh; the greatest fixed point $\nu X.\varphi$ with $X$ positive; and
graded separation $\varphi \mid_k \psi$, which splits a term under $\mathsf{H}$ with exactly $k$
shared restricted names \cite[\S2.3]{ClassifierNote}.

\begin{remark}[Restriction is sugar]
\label{sci:rem:sugar}
We treat \texttt{new} as syntactic sugar throughout. In a reflective calculus a name is a quoted
process, so a ``fresh'' name is manufactured by a structural scheme rather than drawn from a stock of
atoms, and freshness is relative to a scope rather than global. We remain agnostic about the scheme,
subject to one constraint: it must be \emph{decentralised}, manufacturing names from material a
computation already holds rather than by consulting a registry. That constraint is not an engineering
preference. A central name server would be an oracle sitting across the cut and consultable by every
computation, which contradicts the isolation premise of \S\ref{sci:sec:cut} outright.

Two consequences run through the rest of the note. Names have \emph{structure}, so a name predicate
sees all the way into one and the hidden-name quantifier is descriptive rather than opaque; and names
have \emph{size}, so by Definition~\ref{sci:def:sense} they carry a running cost. Concealment is therefore
never secrecy, only price --- which is why $\delta$ was introduced as a cost parameter and not as a
wall.
\end{remark}

\subsection{The structural predicates are admitted, and they are senses}\label{sci:S3.2}

In \cite{ClassifierNote} the structural predicates are wielded by an analyst standing outside the
term, holding it in hand, free to count its gates and test its decompositions. That is not the
scientist's situation, and admitting the structural predicates without saying where they may be
applied would hand the scientist an oracle and destroy the isolation premise of \S\ref{sci:sec:cut}.

The discipline that saves them is that a sense is a structural predicate evaluated \emph{at the
surface}.

\begin{definition}[Depth]\label{sci:depth}
The \emph{depth} $d(\varphi)$ of a structural formula is the maximal nesting of term-former
connectives and prefixes occurring in it, with $d(\top) = 0$, $d(\mathrm{out}(n)\varphi) =
d(\mathrm{in}(n)\varphi) = 1 + d(\varphi)$, $d(\varphi \mid \psi) = \max(d(\varphi), d(\psi))$,
$d(\mathsf{H}x.\varphi) = 1 + d(\varphi)$, and $d(@\varphi) = d(\varphi)$.
\end{definition}

\begin{definition}[Sensing, and its price]
\label{sci:def:sense}
For $W \equiv S \mid E$, the scientist $S$ may evaluate $\varphi$ against $E$ whenever
$d(\varphi) \le \alpha$, where $\alpha$ is $S$'s \emph{acuity}; the evaluation debits $S$'s stack by
$\kappa_{\mathrm{sense}}(d(\varphi))$, a monotone schedule. Acuity is not a fixed endowment: $\alpha$
is whatever $S$ can currently pay for.
\end{definition}

Two comments. First, the depth grading is what keeps perception from being an oracle: seeing deeper
costs more, so the interior of $E$ is not free, it is merely expensive, and past a point unaffordable.
Second, $d(@\varphi) = d(\varphi)$ rather than $1 + d(\varphi)$ is deliberate and records the third
grade of access: opening a name one holds is not a passage through a surface, because one already
holds it. Reflection is cheap for exactly the reason predation is not.

\begin{remark}[Senses are instruments, and instruments are held]
\label{sci:rem:senses-instruments}
There is an objection to this whole subsection which the pricing schedule does not answer. If the
scientist's perception \emph{is} the structural predicates, and the structural layer sees more than
interaction does, then interaction was never the whole ceiling and \S\ref{sci:sec:cut}'s three grades
were not exhaustive after all. Chapter~\ref{ch:obsext} is the reply, and it is worth stating in this
chapter's vocabulary because it changes what a sense is.

A structural observation can be \emph{represented as} an interaction, by conservatively extending the
theory with destructuring rules: opening a term at a constructor becomes a rendezvous with a probe,
and the existential over decompositions in $\varphi \mid \psi$ becomes an existential over
transitions (Lemma~\ref{ox:lem:exposure}). Perception is therefore not a fourth grade of access
smuggled in beside the three; it is interaction in a theory whose observer holds instruments. But the
instruments are held, not free, and that is the part this chapter had not said. Which structural
predicates a scientist may apply is the parameter $\Ob$ of Definition~\ref{ox:def:ext}, and it is a
fact about the learner and its world, alongside acuity and alongside the budget.

Two consequences. Table~\ref{sci:tab:access}'s first row lists instruments the scientist is here
assumed to hold in full; a learner holding fewer sits at a coarser point of
Proposition~\ref{ox:prop:mono} and has a smaller ideal logic, not merely a smaller affordable
fragment. And $\kappa_{\mathrm{sense}}(d)$ cannot be derived from the construction:
\S\ref{ox:sec:cost} shows that depth bounds the sequential destructuring path but that checking cost
depends on branching and on evaluation strategy besides. The schedule remains a modelling choice,
which \S\ref{sci:S14.2} already concedes; what is new is that we now know why no derivation was
available.
\end{remark}

\begin{table}[t]
\centering
\fitwidth{%
\begin{tabular}{@{}llll@{}}
\toprule
\textbf{grade} & \textbf{instrument} & \textbf{applies to} & \textbf{price} \\
\midrule
perception & $\mathrm{out}, \mathrm{in}, \mid, \mid_k, \mathsf{H}$ & the surface of $E$ &
$\kappa_{\mathrm{sense}}(d)$, monotone in depth \\
interaction & $\langle K_j \rangle \varphi$ & the behaviour of $E$ & one rendezvous per step, and time \\
reflection & $@\varphi$ & names $S$ holds & a computation on a name; no communication \\
\bottomrule
\end{tabular}}
\caption{The three grades of access. The ladder is also an ordering by invasiveness
(\S\ref{sci:sec:exposure}), and in the opposite direction.}
\label{sci:tab:access}
\end{table}

\section{Hypotheses, assays, and the \texttt{where} clause}\label{sci:S4}
\label{sci:sec:assay}

\subsection{The embedding, in operational position}\label{sci:S4.1}

An embedding of the logic into the calculus is what lets \rhoc\ computations compute over hypotheses.
The obvious construction --- reify formulae as data, write an interpreter --- works but is
unsatisfying, because the scientist must then carry the interpreter and the interpreter is not part of
the theory. The variant of \rhoc\ whose for-comprehension carries a \texttt{where} clause
\cite[\S9]{Omnibus} does better. Its general form is
\[
\texttt{for}\big( p_{11} \leftarrow c_{11} \,\&\, \cdots \,;\, \cdots \,;\, p_{1n} \leftarrow c_{1n}
\,\&\, \cdots \ \texttt{where}\ \vartheta \big) P
\]
with $\&$ joining within a group, $;$ separating groups, and $\vartheta$ a condition drawn from the
booleans over ground formulae, spatial connectives, and the context-labelled modality, in which the
variables bound by the preceding clauses may be mentioned. A guard is a formula in operational
position. Hence:

\begin{quote}
\itshape A hypothesis installed as a guard is an experiment, and a confirmation is a rendezvous.
\end{quote}

The checker is the reduction rule. Nothing is carried, nothing is interpreted, and the cost of testing
a hypothesis is metered by the same apparatus that meters everything else --- which is essential here,
since an unpriced hypothesis test would be an oracle and would collapse the entire motivational story
of \S\ref{sci:sec:game}.

\subsection{The assay}\label{sci:S4.2}

\begin{definition}[Assay]
\label{sci:def:assay}
An \emph{assay} is a pair $(K, \varphi)$ of a probe context $K$ and a predicted property $\varphi$,
realised on a probe channel $p$ as
\begin{lstlisting}
   K[ ... ]                                        // the probe: perturb the environment
 | for (x <- p where phi(x)) { confirm!(*x) }       // prediction holds
 | for (x <- p where ~phi(x)) { refute!(*x) }       // prediction fails
\end{lstlisting}
\end{definition}

The negation in the second guard is what makes this more than an expression of hope. In a concurrent
setting a guard that simply fails to fire is indistinguishable from a guard still waiting, so a naive
encoding of falsification cannot tell refutation from patience. The complementary pair repairs this.

\begin{proposition}[Trichotomy]
\label{sci:prop:trichotomy}
Let $A = (K,\varphi)$ be an assay on probe channel $p$, with $\varphi$ decidable on the received term
within the ambient budget. Then for each message arriving on $p$ exactly one of the two guards is
enabled, and the assay yields one of three outcomes: \textsc{confirm}, \textsc{refute}, or $\bot$,
where $\bot$ occurs exactly when no message arrives on $p$ before the budget for $A$ is exhausted.
\end{proposition}

\begin{proof}[Proof sketch]
$\varphi$ and $\lnot\varphi$ partition the terms on which both are decidable, so at most one guard is
enabled per message and the two receipts do not race. If a message arrives and the budget suffices to
decide $\varphi$ on it, one guard fires. Otherwise no output occurs on either \texttt{confirm} or
\texttt{refute}, which is the residual case.
\end{proof}

We stress that $\bot$ is not a defect. It is the honest residue: the environment did not answer, and
no amount of logical apparatus can distinguish that from an environment that would have answered
later. What it does is separate that genuine ignorance from falsity, which the naive encoding conflates.

\subsection{Refutation is budget-relative, and that fixes the fragment}\label{sci:S4.3}

\begin{proposition}[Asymmetry of the modalities]
\label{sci:prop:asymmetry}
Two distinct asymmetries hold, and it is worth separating them because the note uses both.
\begin{enumerate}
\item \emph{Over positions.} Let $\varphi = \exists j.\,\langle K_j \rangle \psi$. Then
\textsc{confirm} is a finite witness for $\varphi$ --- exhibit the position and the step --- while
refuting it requires ruling out every position, and for $\varphi = \forall j.\,[K_j]\psi$ the
polarity reverses: a single position is a finite refutation and no finite run confirms.
\item \emph{Over time.} For a fixed position $j$, $\langle K_j\rangle\psi$ is confirmed the moment
the step fires, but its negation is not finitely witnessed at all, because a step that has not yet
become enabled is indistinguishable from one that never will. This is Proposition~\ref{sci:prop:trichotomy}'s
$\bot$ seen from the other side.
\end{enumerate}
\end{proposition}

\begin{remark}[Which asymmetry does which job]
\label{sci:rem:whichasym}
The first clause is what bounds the testable fragment, since nesting $\exists j$ under $\forall j'$
is what modal depth counts, and it is the clause invoked immediately below and again in
\S\ref{sci:sec:repro}. The second is what makes the complementary guard pair of
Definition~\ref{sci:def:assay} necessary rather than fussy: without the negated guard, clause 2 says a
guard that has not fired is uninterpretable. Stating the proposition with an uninstantiated $K$
conflates the two, and the conflation is easy to miss because both conclusions are true.
\end{remark}

The consequence is the note's first structural fact about mortality. A scientist with a finite budget
can only run assays whose outcome is decided within that budget, so the fragment of the ideal logic it
can actually test is the one bounded in modal depth and fixed-point unfolding. That fragment is not
adequate --- a logic whose model checking terminates never is, for a Turing-complete host
\cite[\S1.3]{ClassifierNote}. Therefore:

\begin{quote}
\itshape A mortal scientist necessarily holds a non-adequate theory, and which non-adequate fragment
it holds is fixed by what it can pay for. The budget is the inductive bias.
\end{quote}

This is not a limitation we regret. It is the formal content of the observation that finite creatures
have finite theories, and it locates the inductive bias somewhere unusually concrete: not in a prior,
not in a regulariser, but in a token stack.

\section{The hypothesis space is the relaxation lattice}\label{sci:S5}
\label{sci:sec:lattice}

We do not need to invent a space of hypotheses. \cite[\S3.3]{ClassifierNote} constructs one and calls
it something else.

Recall the construction. For a process $P$, a \emph{characteristic formula} $\chi_P$ is satisfied by
exactly the processes equivalent to $P$; it turns an equivalence question into a model-checking
question. Beneath $\chi_P$ sits a lattice of progressively weaker formulae obtained by forgetting
parts of it --- the recursion, the continuations, a private name --- each forgetting yielding a
coarser predicate that more processes satisfy. In \cite{ClassifierNote} this lattice is where a
classifier's dials live. Here it is the space the scientist moves in.

\begin{definition}[Revision]\label{sci:revision}
A \emph{revision} is a move on the relaxation lattice: \emph{generalisation} forgets a conjunct and
moves down, \emph{specialisation} restores one and moves up. A scientist's trajectory is a walk
$\varphi_0, \varphi_1, \varphi_2, \ldots$ on this lattice.
\end{definition}

\begin{definition}[Stratified distance]\label{sci:stratified-distance}
For $\varphi, \psi$ agreeing on all approximants below stage $n$ and differing at $n$, set
$d(\varphi, \psi) = 2^{-n}$, where the stages are the $\nu$-unfolding approximants
$\nu^0 \supseteq \nu^1 \supseteq \cdots$. This is an ultrametric.
\end{definition}

The stratification is not chosen for convenience: it is how bisimulation is already approximated, it
grades with modal depth, and by Proposition~\ref{sci:prop:asymmetry} modal depth is what the budget
bounds. Distance, depth, and cost are therefore one graded structure rather than three, which is the
main reason to prefer this metric to a measure-theoretic one on denotations.

\begin{proposition}[The limit of inquiry is unaffordable]
\label{sci:prop:limit}
Let $E$ be an environment whose behaviour is not eventually structural in the sense of
\cite[Rem.~6]{ClassifierNote} --- that is, whose recursion follows the reduction rather than a finite
structural cycle. Then $\chi_E$ is not exact at any finite unfolding, and no scientist with a finite
token supply can install it as a guard.
\end{proposition}

So $\chi_E$ --- the complete theory of the environment --- sits at the top of the lattice as a limit
the scientist approaches and never reaches. Science converges \emph{along} the lattice.

\begin{remark}[When a finite theory is complete]
\label{sci:rem:structural}
The converse case is worth naming because it is the good one. \cite[Rem.~6]{ClassifierNote} observes
that a recursion following the \emph{structure} rather than the reduction is exact at a finite
unfolding: for a diagram with $2n$ arcs, $\nu^{2n} = \nu$ on the nose. Read for the scientist, this
says that a finite theory is complete exactly when the environment's regularity is structural.
Recursions that follow the reduction are the unbounded ones. This is a serviceable formal criterion
for which subject matters admit closed sciences and which do not, and we suspect it is more useful
than it looks.
\end{remark}

\section{Searching the lattice}\label{sci:S6}
\label{sci:sec:search}

Section~\ref{sci:sec:lattice} supplies a space of hypotheses and a metric on it, and says that a scientist's
trajectory is a walk. It does not say how the walk is chosen. That is the revision policy, which we
left open as one of the note's three free parameters (\S\ref{sci:sec:audit}), and we leave it open still.
What we do here is narrower and, we think, more useful: the metric is an \emph{ultra}metric, and
ultrametric spaces are geometrically unlike the spaces search algorithms are usually designed for. Much
of what one would reflexively try is unavailable, and one thing one might not have thought necessary
turns out to be forced.

\subsection{The geometry is a tree}\label{sci:S6.1}

Three facts about ultrametric spaces do the work. Balls are nested or disjoint, never partially
overlapping. Every point of a ball is a centre of it, so no point in a neighbourhood is distinguished.
And there is no betweenness: for $\varphi \ne \psi$ there is no formula standing partway from one to
the other in the sense a convex combination would.

\begin{remark}[There is no gradient]
\label{sci:rem:nogradient}
Gradient methods are not merely a poor fit here; they are not well typed. A gradient presupposes a
direction in which one may move by an arbitrarily small amount and a notion of interpolation between
neighbouring points, and an ultrametric space supplies neither. What it supplies instead is branching,
so the search available on a hypothesis lattice is tree search, and the strategies below are all tree
strategies.
\end{remark}

This has a consequence which we regard as the central fact governing revision policies.

\begin{proposition}[Confinement]
\label{sci:prop:confinement}
Let $\varphi_0, \varphi_1, \ldots, \varphi_k$ be a revision walk with
$d(\varphi_{i-1}, \varphi_i) < 2^{-n}$ for every $i$. Then $d(\varphi_0, \varphi_k) < 2^{-n}$: the
walk never leaves the ball of radius $2^{-n}$ in which it began.
\end{proposition}

\begin{proof}
By the ultrametric inequality $d(\varphi_0, \varphi_k) \le \max_i d(\varphi_{i-1}, \varphi_i)$, and
the maximum is below $2^{-n}$ by hypothesis.
\end{proof}

In a Euclidean space, sufficiently many small steps reach anywhere; here they reach nothing outside the
ball they started in, however many are taken and however long the scientist lives. A policy of
incremental revision is not slow at escaping a bad basin --- it cannot escape one at all. Whatever ball
a scientist's early commitments place it in is the ball it will die in, unless it makes a move that is
not small.

We take this up again in \S\ref{sci:sec:repro}, because it turns a biological mechanism into a
computational necessity.

\subsection{The two parameters}\label{sci:S6.2}

\begin{definition}[Revision move]
\label{sci:def:move}
A \emph{revision move} is a pair: a \emph{locus} in the formula tree --- which structural predicate is
to be varied --- together with an \emph{operation} at that locus, weakening or strengthening. The
locus fixes the direction; its depth fixes the distance.
\end{definition}

So the two parameters the policy must choose are exactly the two the metric grades.

\paragraph{Distance is backtracking depth.} A move of distance $2^{-n}$ agrees with the current
hypothesis on every approximant below stage $n$ and differs at $n$. Small distance means a deep
backtrack and a fine adjustment; large distance means a shallow backtrack and the abandonment of an
early commitment.

\paragraph{Direction is a choice of sibling.} At a node of the tree the available directions are the
one-step refinements the generated logic admits at that locus. The branching factor is therefore fixed
by which connectives OSLF emitted, and is not ours to choose.

The price schedule follows without further assumption.

\begin{proposition}[Cost is monotone in distance]
\label{sci:prop:movecost}
Evidence gathered by an assay of modal depth $k$ bears on the approximants at stage $k$ and above. A
move of distance $2^{-n}$ therefore invalidates exactly those confirmations obtained at depth $n$ or
greater, and the count of invalidated confirmations is non-decreasing as the distance of the move
increases.
\end{proposition}

A radical revision is expensive in precisely the currency the scientist has been spending: it discards
the assays that bought the commitments it abandons. Cheap moves are confined
(Proposition~\ref{sci:prop:confinement}) and unconfined moves are dear, which is the trade-off any policy
is negotiating.

\begin{remark}[Two gradings, one poset]
\label{sci:rem:twogradings}
The relaxation lattice and the ultrametric grade the same partial order differently, and conflating
them is easy. The lattice orders by logical strength: how much of a characteristic formula has been
forgotten. The metric orders by depth: how far down the forgotten material sat. Forgetting a conjunct
at stage $1$ and forgetting one at stage $5$ are both a single lattice step, and are metrically very
far apart. Definition~\ref{sci:def:move} keeps the two apart by naming the locus and the operation
separately.
\end{remark}

\subsection{Strategies}\label{sci:S6.3}

We catalogue rather than recommend, since which policy is best is an ecological question and depends on
the distribution of \S\ref{sci:sec:distributions}.

\begin{description}
\item[Conservative.] Always move the least distance that the last refutation forces. Cheap by
Proposition~\ref{sci:prop:movecost}, low variance, and by Proposition~\ref{sci:prop:confinement} permanently
confined to its initial ball.

\item[Radical.] On refutation, abandon a commitment high in the tree. Unconfined, and expensive in
exactly the confirmations it throws away.

\item[Adaptive, with refutation rate as temperature.] Jump far when refutations are frequent --- a
high refutation rate is evidence that a commitment high in the tree is wrong --- and refine when they
are rare. This is annealing-shaped, but the temperature is not a schedule imposed from outside: it is
the ratio of \texttt{refute} to \texttt{confirm} events, which the assays of \S\ref{sci:sec:assay} already
produce as messages. We regard this as the default policy a designer would reach for, precisely
because its control parameter is already an observable of the system.
\end{description}

Direction selection is the second half of the policy, and the candidates differ in what they optimise.

\begin{description}
\item[Blind.] Choose a sibling uniformly. Requires no evaluation, and is what a lineage does when the
choice function of Definition~\ref{sci:def:crossover} is unbiased.

\item[By formula cost.] Prefer siblings whose formulae are cheap to check. This is the minimum
description length bias of \S\ref{sci:sec:scores} appearing as a search heuristic rather than as an
accounting identity, and it has the merit that the scientist prefers hypotheses it can afford to test.

\item[By discriminating power.] Prefer a sibling $\psi$ for which some assay separates $\psi$ from the
current $\varphi$ --- that is, design the crucial experiment first and let it select the direction.
This is the distinguishing-formula problem for bisimulation inequivalence \cite{Cleaveland}, and it is
the most classically scientific of the options: it maximises the information a single assay returns.

\item[By expected yield.] Prefer siblings which, if true, locate more tokens.
\end{description}

\begin{remark}[Yield, not truth]
\label{sci:rem:yieldbias}
The last of these is the one this world actually selects for, and it is worth saying so plainly. A
mortal scientist is not rewarded for holding true hypotheses but for holding nutritive ones, so a
lineage under selection will drift toward yield-biased direction choice and away from
discrimination-biased choice wherever the two disagree. This is the pragmatist objection of
Remark~\ref{sci:rem:pragmatist} arriving at a specific mechanism rather than as a general caveat: it is not
that the scientist is indifferent to truth, but that its search is steered at every node by what pays.
\end{remark}

\subsection{The search is itself metered}\label{sci:S6.4}

One further constraint is easy to overlook and is not a detail. Choosing a direction requires
evaluating candidate siblings, and evaluation is model checking, and model checking is priced. A
scientist therefore cannot survey the branching it faces.

\begin{proposition}[Bounded branching]
\label{sci:prop:branching}
The number of siblings a scientist can evaluate before committing to a move is bounded by its stack
divided by the cost of checking one. Hence the effective branching factor of its search is a function
of its wealth, not of the logic.
\end{proposition}

\begin{corollary}[Poverty is confining]
\label{sci:cor:poverty}
A scientist whose stack affords few evaluations per move is likelier to take the cheapest available
move, and by Proposition~\ref{sci:prop:movecost} the cheapest moves are the shortest. By
Proposition~\ref{sci:prop:confinement} a policy of short moves cannot leave its ball. Impoverishment
therefore tends to confinement, and the confinement is not recoverable by patience --- only by a
subsequent move that the scientist could not afford to evaluate.
\end{corollary}

We note without developing it that this closes a loop with \S\ref{sci:sec:trophic}: a source-feeder whose
inquiry converges accumulates the surplus that would buy a wide survey, and has nothing left to search
for; a prey-feeder that needs the survey is the one paying an unbounded epistemic rent.

\subsection{What escapes}\label{sci:S6.5}
\label{sci:sec:escape}

Proposition~\ref{sci:prop:confinement} leaves an individual scientist with exactly one route out of a bad
basin: a move that is not small, at a cost given by Proposition~\ref{sci:prop:movecost}, chosen from a
branching it cannot afford to survey. That is a poor prospect, and it is the reason
\S\ref{sci:sec:repro} is not an appendix to this note but a part of its argument. A crossover at a locus
high in the term tree is, transported along the characteristic-formula construction, a shallow
re-branch in the formula tree --- a displacement of exactly the kind an incremental walk provably
cannot make, and one no individual pays for, because the cost is borne by the offspring that does not
yet exist. The escape is not unconditional: by Remark~\ref{sci:rem:dial} how far a swap can displace is
fixed by the recombination discipline, and the disciplines that displace furthest are the ones that
kill most offspring.

\section{The game}\label{sci:S7}
\label{sci:sec:game}

\subsection{Conservation}\label{sci:S7.1}

Tokens are conserved and never minted. Write $\Theta$ for the total supply of the world. The cost
monad's ledger law and the history monad's record together give, globally,
\[
\sigma + \kappa \;=\; \sigma_0 \;=\; \Theta ,
\]
where $\sigma$ is the total still held in live stacks and $\kappa$ the total spent. Spent tokens do
not vanish; they migrate into the record. So the free supply is monotonically non-increasing, the
accumulation of $\kappa$ is the arrow of time, and there is no primary producer. The whole game is a
finite race.

\begin{remark}[Science is entropy production]\label{sci:science-entropy-production}
The aggregate effect of every scientist in the world is to convert free tokens into history. This is
the precise form of the observation that the cost of remembering is the cost of doing science: the
record is made of spent tokens. It also means predation is not profit in any global sense --- it is
the purchase of someone else's remaining time.
\end{remark}

\subsection{Working in the image}\label{sci:S7.2}
\label{sci:sec:image}

Because \rhoc\ is Turing complete, cost-accounted \rhoc\ translates back into \rhoc: there is an
internalising encoding $\enc{-} : \Crho \to \rhoc$ that represents the located token stacks in the
undecorated calculus \cite[\S6]{Omnibus}. Without loss of generality we may therefore avail ourselves
of the cost-accounted syntax where that is convenient and of the channels at which token stacks are
located where that is convenient.

\begin{definition}[Metabolic channel]
\label{sci:def:metabolic}
Under $\enc{-}$, a located stack becomes a message at a name. We call that name the \emph{metabolic
channel} $m_P$ of the computation $P$, and we call
\[
m_P!(\sigma) \;\mid\; \texttt{for}(t \leftarrow m_P)\{\, \text{debit};\ m_P!(\sigma')\ \mid\ \cdots \,\}
\]
the \emph{reclaim--debit--reissue} cycle. A mortal computation is a term equipped with a metabolic
channel executing this cycle, with a strictly positive burn rate.
\end{definition}

\begin{remark}[Metabolism is a race condition]
\label{sci:rem:race}
Living requires exposing one's stack as a message, if only for the instant between reclaiming and
reissuing. Computing is therefore racing others to one's own stack. We did not put this in; it is
what the comm rule does with Definition~\ref{sci:def:metabolic}.
\end{remark}

\begin{definition}[Harvest, death]
\label{sci:def:harvest}
A \emph{harvest} of $P$ by $S$ is the receipt $\texttt{for}(t \leftarrow m_P)$ performed by $S$. It is
an ordinary communication and requires no new rewrite rule. $P$ \emph{dies} when it cannot fund its
next step: the linear receipt consumed the message it was going to reclaim.
\end{definition}

So eating is a receive, death is starvation, and defence is restriction. All three are already in the
calculus.

\begin{remark}[One ledger]\label{sci:one-ledger}
The translation is not cost-free: simulating the metering costs \rhoc-steps. We therefore adopt one
metering convention throughout and state it once. \emph{Costs are always read in the decorated
calculus $\Crho$; the image is used only for the mechanics of transfer.} This matters because the
central inequalities of \S\ref{sci:sec:exposure} and \S\ref{sci:sec:distributions} are statements about ratios
of costs, and a non-uniform translation overhead would move their boundaries. It is the same
discipline of keeping levels distinct that governs running $\mathcal{C}$ on an
$\mathcal{H}$-equipped model.
\end{remark}

\subsection{Predation is the failure of internalisation to be faithful}\label{sci:S7.3}

Here is the observation that we take to be the note's chief structural contribution.

In $\Crho$ a located stack is inviolable. That is not an accident of presentation; it is precisely
what the located-stack discipline and the ambient-authority fix were built to guarantee. No rewrite
takes your stack, because the stack is not the sort of thing a rewrite addresses.

In the image $\enc{\Crho} \subseteq \rhoc$ the stack is a message at a name, and any process holding
that name may receive it.

These are not the same situation, and the encoding is not faithful against arbitrary \rhoc\ contexts.
It is faithful against \emph{images of $\Crho$ contexts} --- which is exactly the refinement of
\cite{Probing}: an encoding may only be probed by the encodings of source contexts, since arbitrary
target contexts see strictly more and would never permit bisimulation preservation. Therefore:

\begin{theorem}[Predators are non-image contexts]
\label{sci:thm:predator}
Let $\enc{-} : \Crho \to \rhoc$ be the internalising encoding and let $C[-]$ be a \rhoc\ context. If
$C$ lies in the image of $\enc{-}$ on contexts, then $C$ cannot harvest: the stack of any
$\enc{P}$ plugged into $C$ is preserved, since the source discipline forbids the corresponding step.
If $C$ lies outside the image and holds $m_P$, it may harvest. Hence the predatory contexts are
exactly the non-image contexts holding a metabolic name.
\end{theorem}

\begin{proof}[Proof sketch]
Faithfulness of $\enc{-}$ against image contexts is the hosting condition of \cite{Probing}: every
step available to $C[\enc{P}]$ with $C = \enc{C'}$ is the image of a step of $C'[P]$ in $\Crho$, and
no step of $\Crho$ removes a located stack from a term. A context outside the image is under no such
constraint; a bare receipt on $m_P$ is such a context and takes the stack.
\end{proof}

\begin{quote}
\itshape The decorated calculus is the world in which property rights hold. Its image is the state of
nature. The ecology lives in the gap, and eating is what the erasure theorem's side condition was
protecting against.
\end{quote}

This gives the game a rulebook rather than a stipulation.

\begin{definition}[The game]
\label{sci:def:game}
A \emph{game} is a choice of admitted context class $\mathcal{A}$ with
$\enc{\Crho\text{-}\mathrm{Ctx}} \subseteq \mathcal{A} \subseteq \rhoc\text{-}\mathrm{Ctx}$.
\end{definition}

The two endpoints are degenerate: $\mathcal{A} = \enc{\Crho\text{-}\mathrm{Ctx}}$ is decorated
cost-accounted \rhoc\ with no game at all, and $\mathcal{A} = \rhoc\text{-}\mathrm{Ctx}$ is total war
in which any context may take anything it can address. The interesting point is the middle one:

\begin{definition}[Capability gating]\label{sci:capability-gating}
The \emph{capability-gated} game admits those contexts that \emph{derive} $m_P$ from structure they
can perceive, and excludes those that are simply handed it.
\end{definition}

Since names are unforgeable, derivation is real work, and it is the work namespace logic does:
describing a namespace by a property rather than enumerating it. Under capability gating, foraging is
name discovery and a hypothesis is a namespace description that either does or does not locate real
tokens. Everything below is stated for the capability-gated game.

\subsection{The foraging inequality}\label{sci:S7.4}

Conservation has an immediate consequence which we will lean on repeatedly: harvest is not
automatically profitable, because finding and opening a stack costs tokens drawn from the same
conserved supply.

\begin{proposition}[Foraging inequality]
\label{sci:prop:foraging}
A harvest of $P$ by $S$ credits $S$ on balance only if
\[
\sigma_P \;>\; \kappa_{\mathrm{sense}}(\delta_P) \;+\; \kappa_{\mathrm{assay}} \;+\;
\kappa_{\mathrm{break}},
\]
where $\sigma_P$ is the stack recovered and $\delta_P$ the depth at which $m_P$ lies concealed.
\end{proposition}

\begin{corollary}[The profitable band]
\label{sci:cor:band}
Since a well-stocked computation can afford deeper concealment --- raising $\kappa_{\mathrm{sense}}$
--- while a poorly stocked one offers little $\sigma_P$, viable prey occupy a band: too poor is not
worth opening, too rich is not worth cracking.
\end{corollary}

This is optimal foraging theory, and we did not import it; it is what Proposition~\ref{sci:prop:foraging}
says once concealment is priced by Definition~\ref{sci:def:sense}. It is also the inequality whose sign
change draws the phase boundaries of \S\ref{sci:sec:distributions}.

\section{Destructive assay forces a logic of namespaces}\label{sci:S8}
\label{sci:sec:forced}

Harvest ends the prey's computational life. That single decision has a consequence for the hypothesis
language which we did not expect and which we now regard as the best available answer to ``why
namespace logic.''

\begin{proposition}[One specimen, one measurement]
\label{sci:prop:onespecimen}
Let $P$ be a mortal computation and let $S$ harvest it. Then no further assay may be run against $P$.
Consequently a hypothesis whose subject is the individual $P$ has zero expected yield after its own
first test.
\end{proposition}

\begin{proof}[Proof sketch]
By Definition~\ref{sci:def:harvest} the harvest consumes $m_P$'s message and $P$ cannot fund a further
step, so $P$ contributes no further transitions and no further surface. A hypothesis $\varphi$ whose
denotation is $\{P\}$ is therefore, after the test, a predicate with empty extension among live
computations; whatever it cost to form and hold, it cannot be recouped.
\end{proof}

\begin{corollary}[Induction is forced]
\label{sci:cor:induction}
In an ecology with destructive assay, the only hypotheses with positive expected yield are those whose
subject is a \emph{class} of computations. The hypothesis language must therefore be able to describe
a collection by a property rather than by enumeration.
\end{corollary}

That is namespace logic's definition. So the logic is not adopted for elegance or for uniformity with
a companion note; it is what survives the game. Induction --- the move from the specimen to the kind
--- is not an epistemic virtue the scientist is asked to display but the only strategy that pays when
each specimen affords exactly one measurement.

\begin{remark}[The biological reading]\label{sci:biological-reading}
Destructive assay is not exotic. Large parts of experimental biology sacrifice the animal, and the
inferential structure of those fields is exactly the one Corollary~\ref{sci:cor:induction} describes:
because you cannot re-measure the individual, everything must be said about the strain, the cohort, or
the genotype. The formal setting reproduces the methodological consequence, which we take as evidence
that the modelling is not merely decorative.
\end{remark}

\section{Exposure, and the price of being known}\label{sci:S9}
\label{sci:sec:exposure}

\subsection{The exposure lemma}\label{sci:S9.1}

The prey's optimal defence in a world of predators looks, at first, unbeatable. Restrict everything.
A fully closed computation is a live $\tau$-network with no free names; no context can communicate
with it; and by the pessimism of \cite[\S1]{ClassifierNote} its observable behaviour carries almost no
information. It cannot be perceived, so its metabolic channel cannot be derived, so it cannot be
harvested.

It also cannot eat.

\begin{lemma}[Exposure]
\label{sci:lem:exposure}
Let $P$ be a mortal computation with strictly positive burn rate and no free names. Then $P$ performs
no harvest, its stack is monotonically decreasing, and $P$ dies in finitely many steps.
\end{lemma}

\begin{proof}[Proof sketch]
Harvest is a receipt on a metabolic name $m_Q$ belonging to some $Q$ outside $P$
(Definition~\ref{sci:def:harvest}). By Remark~\ref{sci:rem:sugar} there is no primitive restriction to appeal
to, so the hypothesis must be read structurally: $P$ is closed when no name occurring in $P$ is
structurally equal to a name occurring in any other locus. Communication requires such an equality, so
no receipt of $P$ fires against the environment and $P$'s stack receives no credits. With a positive
burn rate it reaches zero in at most $\lceil \sigma_P / \varepsilon \rceil$ steps, where $\varepsilon$
is the minimum debit.
\end{proof}

\begin{remark}[Closure is luck, not construction]
\label{sci:rem:luck}
The structural reading of Lemma~\ref{sci:lem:exposure} is weaker than the restriction-based one in an
instructive way. Because name manufacture is decentralised (Remark~\ref{sci:rem:sugar}), two computations
may independently produce structurally equal names, so a locus cannot guarantee closure by
construction --- only observe that it currently obtains. Perfect isolation is not purchasable. We will
find in \S\ref{sci:sec:repro} that the same fact, read the other way round, is what makes recombination
generative.
\end{remark}

\begin{corollary}[Concealment is a luxury of the starving]
\label{sci:cor:conceal}
Metabolic viability lower-bounds observability: any computation that survives indefinitely must expose
surface, hence must be perceivable, hence is in principle predictable and so in principle edible.
\end{corollary}

The predator's counter-move is now legible, and it is a reading of \cite[Rem.~2]{ClassifierNote}
turned around. Restriction makes a name unobservable but not undescribable; $\mathsf{H}x.\varphi$
reaches under the binder and speaks about what is hidden. A predator therefore attacks what it cannot
observe by \emph{describing} it, which is why the spatial layer and not the modal one is the hunting
instrument, and why senses had to be admitted in \S\ref{sci:sec:logic} for the game to be non-trivial at
all.

\subsection{The invasiveness ladder}\label{sci:S9.2}

Table~\ref{sci:tab:access} ordered the three grades of access by what they reveal. They are ordered in the
opposite direction by what they destroy.

\begin{table}[h]
\centering
\fitwidth{%
\begin{tabular}{@{}llll@{}}
\toprule
\textbf{engagement} & \textbf{effect on the subject} & \textbf{information} & \textbf{nutritive?} \\
\midrule
perceive & none & surface only & no \\
assay & perturbs; the environment changes & behaviour under one probe & no \\
break in & fatal & the whole interior, as structure & yes \\
\bottomrule
\end{tabular}}
\caption{Information gain and subject survival are in strict opposition, and only the fatal end of the
ladder pays.}
\label{sci:tab:ladder}
\end{table}

\begin{proposition}[The opposition is monotone]\label{sci:opposition-monotone}
Along the ladder, the information yielded is non-decreasing and the subject's residual life is
non-increasing; and the only rung that credits the scientist's stack is the last.
\end{proposition}

This is where the note's arc closes. A scientist that is very good at forming hypotheses locates
metabolic channels efficiently; locating them efficiently is eating efficiently; eating efficiently
exhausts the environment that was confirming the hypotheses. \emph{The better the theory, the faster
it destroys the conditions of its own confirmation.} Success is self-terminating, and no additional
assumption was needed to make it so.

\section{Mixed environments}\label{sci:S10}
\label{sci:sec:mixed}

The environment contains a mixture: bare token stacks, non-scientist mortal computations, and other
scientists. Each stratum behaves differently and the third is the interesting one.

\paragraph{Bare stacks} are the primary resource: inanimate, undefended, cheap to open, and finite.
They are what makes the shallow-and-external regime of \S\ref{sci:sec:distributions} a regime in which
science does not pay --- when food lies in the open, senses suffice and a theory is pure overhead.

\paragraph{Non-scientist mortal computations} burn tokens and may have structure, but form no model.
They defend passively, by restriction, and are subject to Lemma~\ref{sci:lem:exposure} like everyone else.

\paragraph{Scientists} model back, and three consequences follow.

\begin{proposition}[Scientists are preferred prey]
\label{sci:prop:preferredprey}
A scientist exposes more surface than a passive mortal computation, since inquiry is a dialogue of
depth at least two conducted over channels that are observable across the cut, whereas metabolism
alone requires only enough surface to harvest (\S\ref{sci:sec:chantypes},
Proposition~\ref{sci:prop:sciencesig}). By Corollary~\ref{sci:cor:conceal} a scientist is therefore more
readily perceived and its metabolic channel more readily derived. The comparison is structural and
needs no normalisation by stack size or burn rate.
\end{proposition}

The most epistemically active agents are the most edible. We do not think this is an artefact.

\paragraph{Theory theft.} Reflection makes another scientist's hypothesis readable. A hypothesis in
guard position is a term; a term may be quoted; and $@\varphi$ inspects the quotation without
communicating. So a scientist may acquire a model by reading rather than by experimenting --- cheaper,
and not fatal to anyone. Whether honest science survives this depends on whether holding a model
requires exposing it, and in this calculus it does, because a guard is surface. We record the
equilibrium question as an open problem (\S\ref{sci:sec:open}) rather than settle it.

\paragraph{Imagination.} Turing completeness plus reflection means a scientist may carry $\enc{-}$
itself as data and \emph{run} a model of a prey's metabolism inside its own budget, paying tokens to
avoid paying tokens.

\begin{proposition}[When imagination pays]
\label{sci:prop:imagination}
Simulating a hunt is viable exactly when the cost of the simulation is below the expected cost of the
errors it prevents:
\[
\kappa_{\mathrm{sim}} \;<\; \Pr[\text{failed hunt}] \cdot \big(\kappa_{\mathrm{sense}} +
\kappa_{\mathrm{assay}} + \kappa_{\mathrm{break}}\big).
\]
\end{proposition}

This is the point at which the scientist stops being a reactive forager and becomes a planner, and it
is where the stacking of the cost and history monads acquires an interpretation that is not merely
formal: the simulated prey is itself a mortal computation, metered inside the predator's budget.

\section{Trophic structure}\label{sci:S11}
\label{sci:sec:trophic}

\subsection{Conservation is the right scope, not a limitation}\label{sci:S11.1}

Smith's organisation of the biosphere \cite{SmithMorowitz} distinguishes organisms that consume other
organisms from those that consume only external energy, and it is tempting to read the second class as
evidence that a model of life needs an open system --- a source outside the conserved pool. We do not
think that reading survives inspection, and we record why, because the temptation is to apologise for
conservation at exactly the point where conservation is doing the work.

The sun is a reservoir of finite extent, situated in a physical system in which energy is conserved.
Its apparent externality is an artefact of where the boundary was drawn: put the boundary at the
biosphere and the sun is outside it; put the boundary at the solar system and there is one conserved
budget being spent down. Openness, here as generally, is a bookkeeping convenience.

So conservation is not a disanalogy with Smith's picture. It is the picture at the correct scope, and
the whole trophic stratification is recoverable without minting anything.

\subsection{The distinction is an access profile}\label{sci:S11.2}

If the sun is a token stack, then the autotroph and the heterotroph are not distinguished by
\emph{what} they eat. They are distinguished by how the stack they eat may be accessed.

\begin{table}[h]
\centering
\fitwidth{%
\begin{tabular}{@{}lll@{}}
\toprule
& \textbf{a source} & \textbf{a prey} \\
\midrule
release & rate-limited; cannot be accelerated & all at once, on demand \\
exhaustible within a lifetime & no & yes \\
exclusive & no; many may harvest at once & yes; a linear receipt \\
adaptive & no & yes, and it models the harvester back \\
\bottomrule
\end{tabular}}
\caption{The trophic distinction is an access profile, not a kind of substance.}
\label{sci:tab:trophic}
\end{table}

The calculus carries this distinction natively, and no new sort is required to express it.

\begin{definition}[Source and prey]
\label{sci:def:source}
A metabolic channel is a \emph{prey} channel when its stack is offered by a linear send, so that a
single receipt exhausts it. It is a \emph{source} channel when its stack is offered by a persistent
emitter --- a replicated or contract-driven send delivering bounded quanta indefinitely --- so that
receipts are non-exclusive and rate-limited by the emitter rather than by the harvester.
\end{definition}

Same sort, same rewrite rule, different multiplicity. Whether a locus is a source or a prey is
therefore a property expressible in the generated logic, and a mortal computation's diet is fixed by
which of these profiles its senses are tuned to derive.

\begin{remark}[Trophic type is a fact about the cut]
\label{sci:rem:cut}
It follows that trophic type is not intrinsic. Whether something counts as external energy or as
another organism depends on which side of the cut of \S\ref{sci:sec:cut} it sits, and at what granularity
one looks: a colony resolves into members, a member into a source and a consumer sharing a boundary.
Since in this setting the boundary is a restriction, ``external energy'' has a formal referent and a
relative one. Smith's stratification is thereby recovered as a statement about how a world is cut, not
as a statement about two kinds of stuff.
\end{remark}

\subsection{The split is predicted, not stipulated}\label{sci:S11.3}

The payoff is that we do not have to posit the trophic division. Two results already in hand imply it.

Remark~\ref{sci:rem:structural} observes that a recursion following the \emph{structure} rather than the
reduction is exact at a finite unfolding. A rate-limited persistent emitter is precisely such a
recursion: a replicated send, structurally regular, non-adaptive. Proposition~\ref{sci:prop:limit}
observes that an environment whose recursion follows the \emph{reduction} admits no characteristic
formula exact at any finite unfolding. An adaptive prey, being itself a mortal computation, is
precisely such an environment.

\begin{proposition}[The two sciences]
\label{sci:prop:twosciences}
A mortal computation feeding on source channels faces an environment whose complete theory is exact at
a finite unfolding: its inquiry converges, and after convergence its epistemic outlay falls to the
rent on a fixed formula. A mortal computation feeding on prey channels faces an environment that
revises against it, so no finite theory is ever complete and a permanent revision cost is incurred for
as long as it lives.
\end{proposition}

\begin{corollary}[Intelligence is a heterotroph's expense]
\label{sci:cor:intelligence}
The standing epistemic cost of a strategy is bounded for source-feeders and unbounded for prey-feeders.
Hence the apparatus of continuing inquiry --- sensing at depth, holding and revising hypotheses,
simulating before acting --- is worth its rent only to the second.
\end{corollary}

Plants do not need brains because the sun is the environment whose complete theory is affordable. We
offer this as a consequence of Proposition~\ref{sci:prop:limit} and Remark~\ref{sci:rem:structural} rather
than as an observation imported from biology; that it is also true of the Earth we take as
corroboration of the modelling rather than as its source.

\subsection{Why the overlap is thin}\label{sci:S11.4}

The space of mortal computations is plainly large enough to contain those that eat both --- nothing in
Definition~\ref{sci:def:source} forbids tuning one's senses to both profiles. On the Earth the
corresponding region is sparsely occupied, and it is worth asking what the formalism says about why.
Four pressures are available, none of them requiring new apparatus.

\begin{enumerate}
\item \textbf{Opposed geometry.} Yield from a rate-limited dilute source scales with collecting
surface; yield from a concentrated one-shot source scales with reach and speed. Surface against
velocity, derived from Table~\ref{sci:tab:trophic} alone rather than borrowed from the shapes of leaves
and predators.

\item \textbf{Opposite sides of Proposition~\ref{sci:prop:limit}.} By Proposition~\ref{sci:prop:twosciences}
the two diets demand different points on the adequacy/checkability trade-off. The source-feeder can
live in a terminating fragment; the prey-feeder cannot, since its subject matter is Turing complete
and adaptive. One budget cannot be allocated to both regimes at once.

\item \textbf{Double rent for single-mode yield.} A foraging strategy is a hypothesis together with a
sensory schedule, and both are held at a cost. A dual-mode computation pays rent on two apparatus
while eating with one at a time --- the argument of \S\ref{sci:sec:scores} applied to strategies rather
than to formulae.

\item \textbf{Hysteresis.} Revision is a walk on the relaxation lattice and walking costs. Where the
transition cost between strategies exceeds the gain from switching, the mixed strategy is not
infeasible but dynamically unstable: it is not a fixed point of the revision dynamics, and
specialisation is an attractor.
\end{enumerate}

The sharpest statement combines the first two with \S\ref{sci:sec:exposure}.

\begin{proposition}[The mixed strategy is the worst of both]
\label{sci:prop:worstofboth}
A computation harvesting both profiles carries the prey-feeder's permanent revision cost
(Proposition~\ref{sci:prop:twosciences}) while maintaining the source-feeder's large collecting surface.
By Lemma~\ref{sci:lem:exposure} and Proposition~\ref{sci:prop:preferredprey} that surface is exactly what makes
it perceivable and its metabolic channel derivable. It therefore bears the heterotroph's epistemic
outlay together with the autotroph's vulnerability.
\end{proposition}

The region is thin because it is dominated, not because it is unreachable.

\subsection{The overlap is populated by composites}\label{sci:S11.5}

There is a further and we think better answer, which is that the question contains a false premise.

On the Earth, dual capability is realised almost entirely by symbiosis: corals with their
zooxanthellae, lichens, kleptoplasty in \emph{Elysia}, and --- decisively --- the mitochondrion and
the chloroplast themselves. Very few single genomes do both at scale; pairs of genomes sharing a
boundary do it constantly.

In this setting that is
\[
\texttt{new } m \texttt{ in } (\,P \mid Q\,) ,
\]
two specialists sharing a metabolic channel under a restriction. Parallel composition makes composites
first-class and graded separation $\mid_k$ measures how much interface they share, so the construction
needs nothing added.

\begin{proposition}[Composite resolution]
\label{sci:prop:composite}
The overlap region is not underpopulated; it is populated by terms that are parallel compositions
rather than atoms. The restriction permits each component to specialise its access profile --- one
tuned to sources, one to prey --- while the shared metabolic channel pools the yield, so
Proposition~\ref{sci:prop:worstofboth} does not apply to the composite.
\end{proposition}

Read with Remark~\ref{sci:rem:cut} this is not a special construction but the general situation: the
organism boundary is drawn by the restriction, not by the strategy, and what looks like one
dual-feeding thing at one granularity is two specialists at another.

\subsection{Succession}\label{sci:S11.6}

One consequence of taking conservation as the correct scope, rather than as a limitation to be
repaired, deserves recording. Because a source is finite --- rate-limited, but not inexhaustible --- a
world does not settle into stable coexisting trophic layers. It runs through them.

\begin{proposition}[Trophic succession]
\label{sci:prop:succession}
While sources deliver at a rate satisfying Proposition~\ref{sci:prop:foraging}, source-feeding dominates
and the cheap fragment of \S\ref{sci:sec:lattice} suffices. As sources deplete, the foraging inequality
fails for them before it fails for prey, since prey concentrate what sources have already collected.
Prey-feeding, and with it the permanent revision cost of
Proposition~\ref{sci:prop:twosciences}, is therefore forced rather than chosen. Thereafter $\kappa$
accumulates until the inequality fails for every locus.
\end{proposition}

The stratification is temporal as much as structural. On a stellar timescale this is a distant
prospect and not a modelling artefact; it is the correct statement of the situation at the scope at
which energy is conserved.

\subsection{An ensemble test}\label{sci:S11.7}

Both claims of this section --- that the mixed strategy is dominated, and that the composite is not
--- are checkable inside the framework rather than merely analogical, and the machinery is that of
\S\ref{sci:sec:distributions}. Sample ecologies across the $(\iota, \delta)$ plane and ask, first, whether
the mixed access profile is ever a viability optimum for an atomic computation, and second, whether the
composite of Proposition~\ref{sci:prop:composite} dominates it wherever it is not. We regard this as a
sharper first computation than calibrating the phase boundaries, because it can come out either way
and one would want to know which.

\section{Reproduction}\label{sci:S12}
\label{sci:sec:repro}

Harvest ends a computational life, so an ecology with no births is a monotone depletion of agents and
nothing in it persists long enough to be a subject matter. This section supplies reproduction. The
mechanism is reflection, and the interest of it is that the genetic apparatus is not modelled but
found: a genome is a quotation, a locus is a context, and crossover is the exchange of parands between
two quoted terms.

\subsection{The naive route, and why sex is worth its price}\label{sci:S12.1}

The cheap route is already in the calculus. Replication, guarded recursion, a persistent receipt, a
contract: each reinstates a term exactly. This is asexual reproduction, and it is strictly cheaper
than the alternative --- no partner, no alignment, no disclosure, no recombination to compute.
Sexual reproduction produces variants rather than copies, so the two are not different apparatus but
the extremes of a single dial: exact reinstatement at one end, recombination at the other. In the
sexual regime there is no exact recursion anywhere in a lineage.

Why pay for the dial to be off zero? The usual answers are available, but this setting supplies a
sharper one, and it comes from the first half of the note.

\begin{proposition}[Sex attacks the predator's hypothesis language]
\label{sci:prop:sexdefence}
By Definition~\ref{sci:def:game} a predator must \emph{derive} its prey's metabolic name, and by
Corollary~\ref{sci:cor:induction} the hypothesis licensing that derivation is a description of a
namespace. Recombination at homologous contexts alters which channels a lineage's members carry.
Hence sexual reproduction perturbs precisely the object about which the predator's hypothesis is
formed, whereas variation confined to continuations perturbs only behaviour under a hypothesis that
remains correctly aimed.
\end{proposition}

This is a Red Queen argument with a mechanism rather than an analogy, and it is available because the
hypothesis language was chosen to be a logic of namespaces in the first place.

\subsection{The germ is a quotation}\label{sci:S12.2}

\begin{definition}[Genome, phenotype]
\label{sci:def:genome}
The \emph{genome} of a mortal computation whose code is $P$ is the name $@P$. Its \emph{phenotype} is
$\ast @P$.
\end{definition}

Everything one wants from the distinction follows from properties the calculus already has.

\begin{itemize}
\item $@P$ does not reduce. A quotation is a name, and names do not step. So a genome does not
metabolise, does not age, and burns no tokens while held.
\item $@P$ is transmissible. Names are what sends carry, so a genome may be handed to a partner or
to an offspring on an ordinary channel.
\item $@P$ is inspectable. By \S\ref{sci:sec:logic} a name predicate sees into a quotation, and by
Remark~\ref{sci:rem:sugar} it sees all the way down. A genome may therefore be screened without being
run --- at reflection prices, which by Table~\ref{sci:tab:access} are the cheapest available.
\item $\ast @P$ costs. Running is what is metered.
\end{itemize}

\begin{remark}[What we do and do not claim here]
\label{sci:rem:noble}
It is tempting to read the first item as a formal Weismann barrier. We decline. What the calculus
gives is an asymmetry of \emph{cost} --- held quotations are free, running terms are metered --- and
that is all we use. Whether an individual's acquired revisions can be written back into what it
transmits is a separate question, on which the empirical case against a strict barrier is
considerable \cite{Noble}, and which is not settled by any theorem here. We note only that a
scientist's hypotheses are held as data rather than as code, which puts them on the near side of any
such barrier, and we leave the matter as Open Problem~\ref{sci:op:noble}.
\end{remark}

\begin{remark}[Persistence and reproduction are distinguished by the ledger]
\label{sci:rem:provision}
Syntactically, a term that reinstates itself and a term that spawns a copy can look alike. Under
conservation they are not alike: a reinstatement that continues to draw on the parent's metabolic
channel is soma, and one that must be endowed from it is offspring. The germ/soma boundary is drawn by
provisioning, not by grammar.
\end{remark}

\subsection{The normal form is a tree, and a locus is a context}\label{sci:S12.3}

Up to structural congruence a process is a parallel composition of primes, and each prime is a
receive, a send, or a drop:
\[
P \;\equiv\; \prod_{i \in I} \texttt{for}(p_i \leftarrow c_i)P_i
\;\Big|\; \prod_{j \in J} c_j!(Q_j)
\;\Big|\; \prod_{k \in K} \ast x_k .
\]
This is a gross decomposition only. Each continuation $P_i$ and each transmitted body $Q_j$ is itself
a process with a normal form of the same shape, so the decomposition recurses and a genome is a
finitely branching \emph{tree} whose nodes are parallel compositions and whose edges are prefixes. The
parands are the individual primes, at every level.

\begin{definition}[Locus, allele]
\label{sci:def:locus}
A \emph{locus} of $P$ is a one-hole context $K[-]$ obtained by deleting a parand from the tree of $P$;
the deleted parand is the \emph{allele} at that locus. A locus is thus identified by the path of
channels from the root to the hole, not by a channel alone.
\end{definition}

The vocabulary is not new: a locus is a one-hole context in the same term algebra the modalities of
\S\ref{sci:sec:logic} are labelled by, so genetics requires no notion of address the note did not
already have. It requires a slightly larger one, and the difference is worth stating rather than
eliding. A modal label $K_j$ is a one-hole context whose hole sits at a \emph{redex}; deleting an
arbitrary parand --- an unmatched send, a receipt with nothing to receive --- yields a one-hole
context that is not the label of any modality. The loci therefore properly contain the modal
labels, and a crossover at a locus outside that class moves the genome to a place no formula of the
generated logic can name.

\begin{proposition}[Loci versus labels]
\label{sci:prop:locivlabels}
Every modal label $K_j$ of $P$ is a locus of $P$; the converse fails. The loci that are labels are
exactly those whose hole is at a communicating pair in the sense of
Definition~\ref{sci:def:intgraph}.
\end{proposition}
\begin{proof}[Proof sketch]
A redex is a parand pair, so deleting it gives a locus, which establishes the inclusion. For the
converse, take $P \equiv c!(Q) \mid R$ with $R$ offering no receipt on $c$: deleting $c!(Q)$ is a
locus, and no rewrite fires there, so no $K_j$ equals it. The characterisation is then immediate
from the definition of a redex position.
\end{proof}

We flag it here and repair it in \S\ref{sci:sec:pairs}, where the repair turns out to be a discipline
we were going to want anyway.

\begin{definition}[Homology]
\label{sci:def:homology}
Loci $K_M$ of the mother and $K_F$ of the father are \emph{homologous} when their channel paths from
the root agree and their binding environments agree --- that is, the prefixes along the two paths bind
the same pattern variables.
\end{definition}

Homology is therefore a partial isomorphism of the two trees rather than a matching of two multisets,
and the binding clause is a reading-frame condition: an allele's body may reference only variables its
new context binds.

\begin{proposition}[Species is a namespace formula]
\label{sci:prop:species}
Two computations can recombine only at their homologous loci, so the extent to which they interbreed
is the extent to which their path sets agree. Reproductive isolation is namespace disjointness, and
conspecificity is expressible as a formula over contexts in the generated logic.
\end{proposition}

\subsection{Crossover}\label{sci:S12.4}

\begin{definition}[Crossover]
\label{sci:def:crossover}
Let $H$ be a set of homologous loci of $@M$ and $@F$ and let $\chi : H \to \{M, F\}$ be a choice
function. The \emph{recombinant} $@C$ is the term obtained by filling each locus in $H$ with the
allele of the parent $\chi$ selects, retaining polarity --- a receive is replaced by a receive, a send
by a send. Where the split is taken between a guard and its body, Definition~\ref{sci:def:homology}'s
binding clause is required and is what keeps the recombinant well formed.
\end{definition}

We take crossover proper rather than assortment of whole primes. Assortment --- exchanging a prefix
together with its body --- is the safe special case in which the binding condition is automatic; it is
Mendel's independent assortment, and it cannot recombine within a gene. Splitting guard from body is
where recombination earns its name.

\begin{proposition}[Effect size is graded by depth]
\label{sci:prop:depthgrade}
A locus at depth $0$ exchanges an entire subsystem; a locus at depth $d$ exchanges detail interior to
a behavioural sequence of length $d$. Hence the perturbation induced by a swap is non-increasing in
the depth of its locus.
\end{proposition}

The gradation is uniform, and it is your own observation about recursion that makes it so: were the
genome to contain recursive definitions, a swap inside a recursive body would be a global edit whose
effect size was large regardless of depth. Since sexual reproduction yields variants and never exact
reinstatement, no such loci occur, and the analogue of the distinction between large-effect regulatory
and small-effect structural variation falls out of the term structure rather than being posited as a
rate.

\subsection{Communicating pairs as units of inheritance}\label{sci:S12.5}
\label{sci:sec:pairs}

Definition~\ref{sci:def:crossover} lets a choice function select any allele at any homologous locus, and
\S\ref{sci:sec:viability} then observes that most of what it produces is inviable. There is a discipline
which removes much of that waste, and it changes what the unit of heredity is.

\begin{definition}[Interaction graph, communicating pair]
\label{sci:def:intgraph}
The \emph{interaction graph} of a genome has the primes as nodes and an edge between a send and a
receive whenever they act on the same channel with compatible pattern. A \emph{communicating pair} is
such an edge together with its two endpoints. The criterion is syntactic by necessity: whether two
primes \emph{will} communicate is a reachability question, and settling it is not something a parent
could afford (Proposition~\ref{sci:prop:selection}).
\end{definition}

\begin{quote}
\itshape The unit of heredity is a redex. A gene is not a piece of structure but an interaction.
\end{quote}

That is worth pausing on, because it obtains something one would otherwise have to impose. Material
that must work together is inherited together --- linkage --- not because a chromosome places it
adjacently but because the pairing \emph{is} the working-together.

It also obtains something we owed. Proposition~\ref{sci:prop:locivlabels} observed that loci properly
contain the modal labels, so an undisciplined crossover can move a genome to an address the logic
cannot name. The discipline closes the gap exactly.

\begin{proposition}[The discipline aligns the two address spaces]
\label{sci:prop:alignment}
Under the communicating-pair discipline, the loci at which crossover may act are precisely the
modal labels $K_j$ of the genome. Hence every crossover is a move between formulae the generated
logic can state, and the transport of \S\ref{sci:sec:transport} is well defined.
\end{proposition}
\begin{proof}[Proof sketch]
Immediate from Proposition~\ref{sci:prop:locivlabels}: the discipline admits exactly the loci whose
hole is at a communicating pair, which is the characterisation of the labels.
\end{proof}

We regard this as the better argument for the discipline, and a different one from waste reduction.
Restricting to communicating pairs does lower the inviable fraction, which is how it was motivated
above. What it also does --- and what nothing else in this note would have supplied --- is make a
crossover and a revision the same kind of object: a change of one-hole context, on one side of the
characteristic-formula map or the other. Without the discipline the two live in different address
spaces and the correspondence is an analogy. With it, it is a bijection.

The discipline does not by itself preserve viability, and the reason is instructive. Suppose the
mother contains a race,
\[
c!(P) \;\mid\; \texttt{for}(x \leftarrow c)Q \;\mid\; \texttt{for}(y \leftarrow c)R ,
\]
and the swap removes the pair $(\,c!(P),\ \texttt{for}(x \leftarrow c)Q\,)$ in favour of a father's
pair on some other channel. The donated interaction is sound; the \emph{residual} receipt on $c$ is
orphaned, which is precisely the deadlock half of the balance condition. Removing a pair cannot break
the material it donates but can break the material it leaves behind.

\begin{proposition}[Balance is preserved exactly by degree-preserving swaps]
\label{sci:prop:degree}
Let a swap remove a communicating pair on channel $c$ and install one on channel $c'$. If $c = c'$,
the number of sends and the number of receipts on every channel are unchanged, no prime is orphaned,
and any race on $c$ survives with the same arity. If $c \ne c'$, the residual primes on $c$ may be
orphaned and the balance condition of \S\ref{sci:sec:viability} may fail.
\end{proposition}

So preserving races is not a second benefit alongside preserving balance; it is the side condition
under which the balance guarantee holds at all.

\begin{remark}[Discipline substitutes for screening]
\label{sci:rem:discipline}
Proposition~\ref{sci:prop:selection} says a parent cannot screen its recombinants. Under a
degree-preserving pair discipline it does not have to screen for balance, because the recombination
operator cannot produce an unbalanced offspring. The cost is paid once, in the machinery that
constrains the operator, rather than per offspring in model checking. This is an argument for why
recombination apparatus should be expected to be \emph{constrained} rather than free, and it depends
on Definition~\ref{sci:def:intgraph}'s criterion being syntactic --- a semantic notion of ``will
communicate'' would hand the saving straight back.
\end{remark}

\begin{remark}[A dial, and it is the dial of \S\ref{sci:sec:search}]
\label{sci:rem:dial}
Three disciplines are now in view: free parand swap, maximal in variation and mostly lethal; pair
swap, which guarantees the donated interaction but risks the residue; and degree-preserving pair swap,
which guarantees both and varies least. Each rung buys viability with explorable distance --- which is
the trade-off of Propositions~\ref{sci:prop:confinement} and \ref{sci:prop:movecost} seen at the scale of a
lineage. It follows that \S\ref{sci:sec:escape}'s claim is conditional: recombination escapes confinement
only if the discipline is loose enough to make a shallow re-branch, and looseness is paid for in dead
offspring. Where the optimum lies is an ecological question, and belongs with the ensemble of
\S\ref{sci:sec:distributions}.
\end{remark}

\begin{remark}[What the discipline does not simplify]
\label{sci:rem:overlap}
Pairs overlap. A send may match several receipts, a receipt's continuation contains further sends
matching elsewhere, and a persistent receipt matches many sends over time, so the interaction graph is
not a matching. ``Swap a pair'' therefore means ``swap an edge together with its endpoints, where the
endpoints carry other edges,'' and Proposition~\ref{sci:prop:degree} is a statement about one channel at a
time. The multi-channel case requires the degrees preserved simultaneously, which is a stronger and
rarer condition, and we do not claim it is generic.
\end{remark}

\subsection{Three kinds of channel, and what a swap can touch}\label{sci:S12.6}
\label{sci:sec:chantypes}

The discipline of \S\ref{sci:sec:pairs} operates on some of a genome and not on all of it, and saying
which requires a classification that is useful well beyond recombination.

\begin{definition}[Channel types]
\label{sci:def:chantypes}
Let $W \equiv P \mid E$. A channel $c$ occurring in $P$ is
\begin{description}
\item[internal] if $P$ contains a matching send and receipt on $c$ and no name structurally equal to
$c$ occurs in $E$: only $\tau$ steps arise from it;
\item[external] if a name structurally equal to $c$ occurs in $E$ but $P$ contains no matching pair on
$c$: its steps are interactions and nothing else;
\item[both] if $P$ contains a matching pair on $c$ \emph{and} a structurally equal name occurs in $E$.
\end{description}
\end{definition}

The classification is relative to the cut, as trophic type was (Remark~\ref{sci:rem:cut}), and for the
same reason: what counts as inside is a question about where the boundary was drawn.

\begin{remark}[This is the surface, made syntactic]
\label{sci:rem:surfacesyntax}
External and both-type channels are exactly the context-observable features of a computation. They
are what a structural predicate evaluated at the surface can reach (\S\ref{sci:sec:logic}), what a
context-labelled modality can exercise, and what Lemma~\ref{sci:lem:exposure} was about: a computation all
of whose channels are internal has no surface, cannot forage, and starves. The three grades of access
in Table~\ref{sci:tab:access} and the three channel types are the same distinction seen from the two sides
of the cut.
\end{remark}

\subsubsection*{Dialogue}

\begin{definition}[Dialogue, and its depth]
\label{sci:def:dialogue}
A \emph{dialogue} is an alternating chain: a step on an external or both-type channel, a causally
connected sequence of internal reductions, and a further step on an external or both-type channel
whose availability depends on the first. Its \emph{depth} is the number of alternations.
\end{definition}

A computation with no dialogue of depth greater than one is a reflex: whatever it emits is not
conditioned on anything but the message immediately received. Depth two and above is conversation ---
the environment's reply is answered in the light of what the first exchange produced. Two consequences
follow, and the first is what the user of this framework will care about.

\begin{proposition}[Conversation is bounded by the stack]
\label{sci:prop:dialoguebound}
Each alternation of a dialogue debits at least the internal chain connecting its two external steps.
Hence the depth of dialogue a computation can sustain is bounded by its token supply, and by
Proposition~\ref{sci:prop:asymmetry} the modal depth at which it can be observed to differ from a rival is
bounded by the same quantity.
\end{proposition}

\begin{proposition}[Inquiry has a structural signature]
\label{sci:prop:sciencesig}
A computation capable of inquiry has a dialogue of depth at least two on both-type channels: the assay
of Definition~\ref{sci:def:assay} emits a probe, receives, conditions on what it received, and emits
again. A computation whose both-type channels support only depth-one dialogue can forage but cannot
revise. Writing the nesting out over redex positions, the signature is
\[
\mathsf{Sci} \;:=\; \exists j.\ \langle K_j \rangle\Big(\exists j'.\ \langle K_{j'} \rangle
\big(\exists j''.\ \langle K_{j''} \rangle \top\big)\Big),
\]
\[
j, j'' \text{ at out-type positions},\qquad j' \text{ at an in-type position},
\]
with all three positions on channels the classification of Definition~\ref{sci:def:chantypes} marks as
both-type. So ``this computation is a scientist'' is a formula of the generated logic rather than a
designation we assign, and with the positions named it is one a model checker could be handed.
\end{proposition}

That repairs something the note has so far taken on trust. The grazer of Appendix~\ref{sci:app:eco} is not
a non-scientist because we declined to give it a hypothesis; it is a non-scientist because its
interaction graph contains no depth-two loop through a both-type channel, and the difference is
checkable.

It also repairs Proposition~\ref{sci:prop:preferredprey}, whose original statement had to normalise by
burn rate. The correct statement is structural: deeper dialogue requires more both-type channels,
and both-type channels are surface.

\subsubsection*{What each kind of swap touches}

\begin{table}[h]
\centering
\fitwidth{%
\begin{tabular}{@{}llll@{}}
\toprule
\textbf{swap} & \textbf{balance} & \textbf{observability} & \textbf{role} \\
\midrule
internal pair & preserved & silent: $\tau$ only & cryptic variation, priced but unseen \\
both-type pair & preserved if degree-preserving & alters the dialogue & adaptive variation \\
external singleton & internally unaffected & alters the interface & a wager on the environment \\
\bottomrule
\end{tabular}}
\caption{The classification says which discipline applies where. Only internal and both-type channels
host pairs at all; external material consists of primes whose counterpart lies in the environment, so
it is swappable singly without orphaning anything --- at the price that its viability is a bet about
the environment rather than a fact about the genome.}
\label{sci:tab:swaptypes}
\end{table}

This grading is independent of the depth grading of Proposition~\ref{sci:prop:depthgrade}. A recombination
policy has two dials: how much structure a swap moves, and how observable the moved structure is.

\begin{proposition}[No neutral variation]
\label{sci:prop:noneutral}
A swap of an internal pair is invisible to weak bisimulation, hence to every context, hence to any
predator's hypothesis. It is not invisible to the ledger: its reductions are metered like any others.
There is therefore no behaviourally silent variation that is also metabolically silent, and selection
acts on internal variation through efficiency rather than through function.
\end{proposition}

\begin{corollary}[Cryptic variation and its release]
\label{sci:cor:cryptic}
Internal variation accumulates without behavioural consequence and becomes consequential exactly when
its channel is reclassified --- which happens when a structurally equal name comes to occur across the
cut. By Proposition~\ref{sci:prop:accidental} recombination can do this, since a transplanted pair may
carry a channel that coincides with a resident external name. A lineage may therefore accumulate
variation cheaply while remaining illegible, and expose it later.
\end{corollary}

This sharpens Proposition~\ref{sci:prop:sexdefence}. Only the observable fraction of recombination attacks
a predator's hypothesis, because a predator's hypothesis is a description of the external and
both-type sub-namespace and nothing else. But internal variation is the reservoir from which
observable variation is later drawn, so the defence is two-staged: accumulate where it cannot be seen,
and expose when the exposure is worth its cost.

\begin{remark}[Dialogue is where the impact is]
\label{sci:rem:impact}
Collecting the section: swapping at redexes concentrates variation on channels that carry
interactions, and among those, the both-type channels are the ones whose interactions are legible
across the cut. A discipline that swaps redexes therefore acts, by construction, on the computation's
dialogue with its environment rather than on structure that no context could ever observe. That is the
sense in which the discipline of \S\ref{sci:sec:pairs} has real consequence: it is not merely a filter for
viability but a filter that concentrates variation where it can matter.
\end{remark}

\subsection{Names, freshness, and the bound on growth}\label{sci:S12.7}

Because \texttt{new} is sugar (Remark~\ref{sci:rem:sugar}), names have size, and size is priced. One might
therefore fear that recombination inflates them: if a child's fresh names were manufactured from its
enclosing term, and its enclosing term combines two parents, names would double with each generation
and a lineage would price itself out of existence within a few dozen.

The calculus forbids it.

\begin{proposition}[Growth is slaved to term size]
\label{sci:prop:namegrowth}
No process contains a name that is its own code, since quotation strictly increases size. Hence every
name occurring in a process is strictly smaller than that process, and no admissible desugaring
manufactures a name from more than a proper part of its scope. Name size cannot drive term growth; it
is bounded by it.
\end{proposition}

The feared scheme is not merely expensive but syntactically impossible, and the bound holds uniformly
across every decentralised scheme --- which is why we may remain agnostic about which one is taken.

What growth remains is the right kind. At a homologous locus the recombinant takes one allele, not
both, so aligned material contributes a pointwise maximum rather than a sum. Growth enters only at
\emph{non-homologous} loci --- material present in one parent alone --- where inheriting the union is
additive per generation. That is gene duplication, it is linear, and by
Definition~\ref{sci:def:sense} it is priced, which is the regime in which selection can act on it.

\begin{proposition}[Decentralisation makes recombination generative]
\label{sci:prop:accidental}
Since names are manufactured locally and freshness is scope-relative, two parents may independently
hold structurally equal names at loci private to each. A recombinant inheriting both then possesses a
single shared channel where its parents had two private ones. Recombination can therefore create
interaction paths neither parent had, rather than only permuting existing ones.
\end{proposition}

Under a central name server, or under a calculus of atomic names with global freshness, this route is
closed. It is the same fact recorded in Remark~\ref{sci:rem:luck}: what denies a locus guaranteed closure
is what allows a lineage novelty.

\begin{remark}[Accumulation, and a reading of senescence]
\label{sci:rem:aging}
Names grow because the terms containing them grow, and a scientist's term grows as it accumulates
probe channels, held hypotheses, and machinery. With sensing priced by depth, every operation then
costs more, and a computation that does not reclaim dead structure eventually cannot afford to run.
This is a discipline-dependent claim rather than a theorem --- garbage collection is the alternative,
and is itself priced --- but under the discipline it says something worth saying: an organism ages
because it accumulates, and accumulation is what learning is.

A condensation pass at birth, canonically renaming the recombinant's private names from a minimal
seed, resets the inherited weight at a cost proportional to genome size. We note that this is an
argument about \emph{cost} and is orthogonal to Remark~\ref{sci:rem:noble}: whatever the traffic between
soma and germ turns out to be, the weight must be renormalised or lineages accumulate without bound.
\end{remark}

\subsection{Viability, and why there is selection rather than design}\label{sci:S12.8}
\label{sci:sec:viability}

A parallel composition of primes is always a term, so syntactic well-formedness is free. Viability is
not, and it is a conjunction of conditions the generated logic can state: closure, the reading-frame
condition of Definition~\ref{sci:def:homology}, balance --- receives without matching sends deadlock,
sends without receipts leak --- and liveness in the marked-graph sense. Screening a recombinant is
therefore model checking, and the classifier apparatus of \cite{ClassifierNote} is the embryology of
this note.

But model checking is metered, and the space is large.

\begin{proposition}[Selection is post hoc because screening is priced]
\label{sci:prop:selection}
Aligned genomes with $n$ homologous loci admit $2^n$ recombinants, where $n$ is the size of the genome
and not the width of its top level. Screening each is an outlay against the same stack that funds
foraging and inquiry. Hence a parent cannot emit only viable offspring --- not for want of a
criterion, but because applying it exhaustively is unaffordable.
\end{proposition}

\begin{quote}
\itshape Reproduction is a gamble because screening is priced. Were model checking free, a parent
would design its offspring and selection would have nothing left to do.
\end{quote}

Neither fine granularity nor cost accounting gives this alone; it needs both.

\subsection{The energy tension: a three-way allocation}\label{sci:S12.9}
\label{sci:sec:energy}

Reproduction takes energy, and under conservation the energy comes from the parents and nowhere else.
An offspring is endowed out of a parent's stack, so every child is a portion of the parent's remaining
life.

\begin{proposition}[Minimum viable endowment]
\label{sci:prop:endowment}
An offspring endowed below the outlay required for its first successful harvest ---
$\kappa_{\mathrm{sense}} + \kappa_{\mathrm{assay}} + \kappa_{\mathrm{break}}$ for the shallowest prey
available to it, by Proposition~\ref{sci:prop:foraging} --- dies before it eats. Provisioning is therefore
bounded below, and the number of offspring a parent can endow is bounded above by its stack divided by
that floor.
\end{proposition}

We take the endowment to be chosen by the parent rather than fixed by the offspring's code size, which
makes the trade-off between many poorly provisioned offspring and few well provisioned ones a strategy
rather than a constant.

The tension the reader will have anticipated is not two-way but three-way. One conserved stack funds
soma maintenance, germ production, and epistemic outlay: the same tokens buy eating, breeding, and
thinking.

\begin{corollary}[Inquiry trades against fecundity]
\label{sci:cor:kselection}
By Corollary~\ref{sci:cor:intelligence} the standing epistemic outlay is unbounded for a prey-feeder. By
Proposition~\ref{sci:prop:endowment} offspring are bounded below in cost. Since both draw on one stack,
a strategy that sustains deep inquiry sustains fewer offspring.
\end{corollary}

This supersedes the reading of surplus offered in \S\ref{sci:sec:scores}: surplus does not simply buy
curiosity, it buys curiosity \emph{or} children, and which it buys is the life-history problem.

\subsection{What reproduction does for the science}\label{sci:S12.10}

Two consequences bear directly on the first half of the note.

\paragraph{The lineage is a second learner.} A scientist's hypotheses are guards; guards are parands;
parands are what crossover exchanges. So the genome contains the hypotheses, and they are revised by
two mechanisms on two timescales --- within a life by walking the relaxation lattice of
\S\ref{sci:sec:lattice}, across lives by recombination. The same object is revised by both. A lineage is
not merely a means of persistence; it is doing science by other means, and the proximate/ultimate
distinction of \S\ref{sci:sec:scores} extends to cover it.

\label{sci:sec:transport}%
The correspondence is tighter than an analogy. A revision move is a locus in the formula tree together
with an operation there (Definition~\ref{sci:def:move}); a crossover is a locus in the term tree together
with a choice of allele (Definition~\ref{sci:def:crossover}). The characteristic-formula construction of
\S\ref{sci:sec:lattice} carries term structure to formula structure, so a swap at a term-context $K_j$
induces a revision at the formula-locus $K_j$ corresponds to --- and by
Proposition~\ref{sci:prop:alignment} the swap is at a $K_j$ in the first place only because the
recombination discipline restricts it to one, without which the induced formula-locus need not
exist. The two mechanisms are one move,
transported along $\chi$, and applied at two timescales. What distinguishes them is the accessible
distance: by Proposition~\ref{sci:prop:confinement} an individual's incremental walk cannot leave its ball,
whereas a swap high in the term tree is a shallow re-branch by construction. \emph{Recombination is how
a lineage makes the displacements its members cannot.}

\paragraph{Reproduction is what makes induction pay.} Corollary~\ref{sci:cor:induction} forced hypotheses
to be about namespaces, since destructive assay leaves an individual-level hypothesis worthless after
its own first test. But a namespace with no live members is worthless too. Reproduction repopulates
the namespace, so a class-level hypothesis retains value across the specimens it costs to test.

\begin{proposition}[Kinds are the condition of inquiry]
\label{sci:prop:kinds}
In an ecology without reproduction the scientist monotonically destroys its own subject matter and no
hypothesis retains extension. With reproduction the population of a namespace is sustained while the
foraging inequality holds for its members, and hypotheses about kinds remain true of live instances.
Species are what make science possible.
\end{proposition}

\begin{remark}[What reproduction does not do]
\label{sci:rem:nofix}
It does not repeal conservation. Tokens are redistributed into new loci and never created, so
Proposition~\ref{sci:prop:succession} stands unaltered and the world still runs down. What changes is the
shape of the decline, not its direction: births convert a monotone extinction of agents into a
population dynamics with a carrying capacity, which is what makes the intermediate band of
Figure~\ref{sci:fig:phase} occupiable for longer than a single generation.
\end{remark}

\section{Distributions over ecologies}\label{sci:S13}
\label{sci:sec:distributions}

\subsection{A world is a partition of a conserved quantity}\label{sci:S13.1}

Because $\Theta$ is fixed and never minted, a world is a partition of $\Theta$ across loci together
with a structural assignment saying which loci are scientists, which are non-scientist mortal
computations, and which are bare stacks. The configuration space is graded by $\Theta$, and a
distribution over it is a microcanonical ensemble in the strict sense --- fixed total, distribution
over microstates. The analogy is exact rather than suggestive, and it is exact because conservation
is exact.

\begin{definition}[Configuration, ecology]\label{sci:configuration-ecology}
A \emph{configuration} is a finite multiset of mortal computations and bare stacks together with an
assignment of stack sizes summing to $\Theta$ and a designation of which computations carry a revision
policy. An \emph{ecology} is a configuration together with a distribution over the configurations
reachable from it.
\end{definition}

\subsection{Two order parameters, and a phase}\label{sci:S13.2}

Two quantities govern the qualitative behaviour and they are independent:
\begin{itemize}
\item $\iota \in [0,1]$, the fraction of $\Theta$ already \emph{internal} to scientists rather than
loose in the environment;
\item $\delta \ge 0$, the mean \emph{concealment depth} of external stacks --- how far under
restriction and structure they lie, which by Definition~\ref{sci:def:sense} sets the sensing cost of
finding them.
\end{itemize}

\begin{figure}[t]
\centering
\fitwidth{%
\begin{tikzpicture}[font=\small, scale=1.0]
  \fill[gray!12] (0,0) rectangle (9,5.2);
  % starvation region (high delta / high iota)
  \fill[pattern=north east lines, pattern color=gray!60, opacity=.55]
    (0,3.5) .. controls (3,3.2) and (5,2.6) .. (6.2,0) -- (9,0) -- (9,5.2) -- (0,5.2) -- cycle;
  % trivial foraging region (low delta)
  \fill[gray!35] (0,0) -- (0,1.3) .. controls (2,1.35) and (3.6,1.1) .. (4.6,0) -- cycle;
  \draw[thick] (0,0) rectangle (9,5.2);
  \node[align=center] at (2.0,0.55) {\itshape foraging\\ \itshape suffices};
  \node[align=center] at (3.1,2.35) {\bfseries science pays};
  \node[align=center] at (6.6,3.9) {\itshape starvation};
  \draw[->,thick] (0,0) -- (9.5,0) node[right] {$\iota$ \ \scriptsize(fraction internal)};
  \draw[->,thick] (0,0) -- (0,5.7) node[above, align=center] {$\delta$\ \scriptsize(concealment depth)};
  \node[below] at (0,-0.1) {\scriptsize $0$};
  \node[below] at (9,-0.1) {\scriptsize $1$};
\end{tikzpicture}}
\caption{The phase diagram. When food lies shallow and in the open (low $\delta$), senses alone
suffice and a theory is overhead: the pure forager beats the scientist. When too much of $\Theta$ is
already internal, or lies too deep to be found for what it yields, discovery costs more than it
returns and everything starves. Hypothesis formation pays in the intervening band. Science is a phase,
not a universal advantage.}
\label{sci:fig:phase}
\end{figure}

\begin{quote}
\itshape Science only pays in an intermediate band --- where food is neither free nor absent.
\end{quote}

We regard this as the most substantive empirical claim the construction makes, and it is a claim the
construction makes rather than one we impose: the boundaries of Figure~\ref{sci:fig:phase} are where the
foraging inequality of Proposition~\ref{sci:prop:foraging} changes sign. The low-$\delta$ boundary is
where $\kappa_{\mathrm{sense}}$ falls below the cost of holding a theory at all; the upper boundary is
where Proposition~\ref{sci:prop:foraging} fails for every available $P$.

\subsection{Sampling}\label{sci:S13.3}

The machinery for turning a configuration into a distribution and sampling trajectories through it is
not new work. Chapter~\ref{ch:weight} takes a continued interactive GSLT, an adequate
spatial--behavioural logic supplying the keys, and a weight map, to a weighted coalgebra
(Definition~\ref{wt:def:coalgebra}) which the Gillespie construction samples
(Theorem~\ref{wt:thm:markov}); assembling the weightings of a fixed theory and taking
the Grothendieck construction gives the fibred form. That is
precisely what an ecology is: the weight map turns a configuration into a distribution, and the
Gillespie construction samples it. \S\ref{sci:sec:functor} takes up what that fibred structure means here.

\section{The audit: what the assumptions already fix}\label{sci:S14}
\label{sci:sec:audit}

It is easy to build a machine that fits a story. We therefore separate, explicitly, what follows from
the commitments and what we chose.

\subsection{Forced}\label{sci:S14.1}

\begin{table}[h]
\centering
\fitwidth{%
\begin{tabular}{@{}p{0.42\textwidth}p{0.5\textwidth}@{}}
\toprule
\textbf{feature} & \textbf{forced by} \\
\midrule
the hypothesis language & OSLF applied to the \rhoc\ presentation; not designed \\
formulae denote bisimulation classes & adequacy of the generated logic \\
exactly three grades of access & the calculus offers no fourth (\S\ref{sci:sec:cut}) \\
senses are structural predicates & admitting the structural layer, priced by depth \\
an assay is a guarded receipt & the \texttt{where} clause \\
refutation is budget-relative & Proposition~\ref{sci:prop:asymmetry} \\
the hypothesis space & the relaxation lattice of \cite{ClassifierNote} \\
hypotheses must be about namespaces & destructive assay (Corollary~\ref{sci:cor:induction}) \\
eating is an ordinary receipt & internalisation plus Turing completeness \\
death is failure to fund the next step & the cost monad \\
$\sigma + \kappa = \sigma_0$ globally & the history monad plus no minting \\
predators are non-image contexts & Theorem~\ref{sci:thm:predator} \\
the exposure lemma & Lemma~\ref{sci:lem:exposure} \\
optimal foraging, and its profitable band & conservation plus depth pricing
(Prop.~\ref{sci:prop:foraging}) \\
information opposes survival & Table~\ref{sci:tab:ladder} \\
a genome is inert, transmissible, inspectable & quotation does not reduce
(Def.~\ref{sci:def:genome}) \\
a locus is a context & the recursive normal form (Def.~\ref{sci:def:locus}) \\
name growth is bounded by term size & no term contains its own code
(Prop.~\ref{sci:prop:namegrowth}) \\
selection is post hoc & screening is priced (Prop.~\ref{sci:prop:selection}) \\
\bottomrule
\end{tabular}}
\end{table}

\subsection{Chosen}\label{sci:S14.2}

\begin{table}[h]
\centering
\fitwidth{%
\begin{tabular}{@{}p{0.42\textwidth}p{0.5\textwidth}@{}}
\toprule
\textbf{parameter} & \textbf{status} \\
\midrule
$\Theta$, the conserved total & free; sets the scale of the whole game \\
the initial configuration and its distribution & free; \S\ref{sci:sec:distributions} \\
the admitted context class $\mathcal{A}$ & free; Definition~\ref{sci:def:game}, and it is the rulebook \\
the pricing schedule $\kappa_{\mathrm{sense}}$, per-rewrite debits & free; the weight map \\
concealment depth $\delta$ of stacks & free; an order parameter \\
harvest all-or-nothing versus partial & free; we take all-or-nothing \\
the revision policy & free; this is the agent design, and we leave it open \\
the desugaring of \texttt{new} & free, subject to decentralisation
(Rem.~\ref{sci:rem:sugar}) \\
provisioning per offspring & free; we make it parent-chosen, so $r/K$ is a strategy \\
the crossover choice function $\chi$ & free; the recombination rate in disguise \\
the observation signature $\Ob$ & free; which term formers the learner may take apart
(Rem.~\ref{sci:rem:senses-instruments}) \\
the map from costs to rates & free; no canonical one exists (Rem.~\ref{sci:prop:section}) \\
\bottomrule
\end{tabular}}
\end{table}

\begin{quote}
\itshape The physics is forced. The ecology and the psychology are the free parameters.
\end{quote}

Everything metabolic and everything epistemic follows from cost-accounted \rhoc\ together with the
logic it generates. What remains to be chosen is the weighting, the initial distribution, and how the
scientist thinks.

We used to regard the smallness of that list as the principal evidence that the modelling choices were
not tuned to the conclusions, and the list has since grown twice. The last two rows are additions of
the present pass, and neither is a detail: the observation signature fixes which distinctions are
available to the learner in principle rather than in budget, and the cost-to-rate map is what an
earlier version of \S\ref{sci:S15.2} claimed the theory supplied for free. A reader auditing this
chapter should assume the list is still short of complete, and should treat the ``forced'' column as a
claim about what follows \emph{given} the choices above rather than as a claim that the choices above
are inconsequential.

\section{Is the assignment functorial?}\label{sci:S15}
\label{sci:sec:functor}

It is natural to ask whether ecologies are a functor of the theory: is there a functor from the
subcategory of cost-accounted graph-structured lambda theories to distributions over computational
ecologies of the kind described above? The answer is instructive, and it is no in a way that tells us
what the right question was.

\subsection{Not from the base}\label{sci:S15.1}

Write $\CGSLT$ for the subcategory of cost-accounted GSLTs. The assignment
$\mathcal{C}(G) \mapsto$ ``its ecologies'' is not determined by $\mathcal{C}(G)$. Two different weight
maps --- two different pricing schedules on the same cost-accounted theory --- give different sensing
costs, hence different foraging inequalities, hence different phase boundaries in
Figure~\ref{sci:fig:phase} and different distributions over trajectories. The object does not determine
the image.

\begin{proposition}[A fibration, not a functor]
\label{sci:prop:fibration}
The ecologies form a fibration $\pi : \Eco \to \CGSLT$ whose fibre over $\mathcal{C}(G)$ is the class
of ecologies presentable in $\mathcal{C}(G)$, indexed by the decorations $\Dec(G)$. The assignment of
distributions is a functor out of the total space --- equivalently, out of the Grothendieck
construction $\textstyle\int_{\CGSLT} \Dec$ --- and not out of the base.
\end{proposition}

This is Chapter~\ref{ch:weight}'s construction with a different label on the target: it already goes
from a Grothendieck construction of weightings into weighted coalgebras, and an ecology-with-a-
distribution is a weighted coalgebra.

\subsection{There is no canonical section}\label{sci:S15.2}

There is nevertheless a rescue. Its shape is the interesting part, and it is smaller than an earlier
version of this section claimed.

\begin{remark}[Cost is not already a weighting]
\label{sci:prop:section}
An earlier version of this section asserted that a cost-accounted GSLT is not decoration-free, on the
grounds that the cost function is itself a weight map; taking costs as the weights would then select a
distinguished decoration for every object of $\CGSLT$, hence a \emph{canonical} section
$s : \CGSLT \to \Eco$ of $\pi$. That is withdrawn. It is not a small correction, because the
canonicity was the whole content of the claim.

Chapter~\ref{ch:weight} is what refutes it, and refutes it in this book rather than from outside.
There a \emph{weight} is a rate: a propensity deciding how often a transition occurs. A \emph{cost}
is a debit against a stack, entering the dynamics through the funding gate $\chi(r,k,\sigma)$, which
decides not how often a transition occurs but \emph{whether it may occur at all}
(Definition~\ref{wt:def:gate}). Those are different data with different types. A debit of five tokens
does not determine a rate of five, nor $1/5$, nor $e^{-5}$; every one of those is a further modelling
choice, and Chapter~\ref{ch:weight} keeps them apart precisely so that a starving agent with a strong
propensity is silent rather than fast.
\end{remark}

What survives is weaker and still useful. Cost accounting canonically supplies the resource ledger; an
ecology additionally requires a pricing or stochastic weighting. Each monotone map from costs to rates
defines a section of $\pi$, and the modeller's choice among them is a real degree of freedom rather
than a departure from the theory's own answer, because the theory does not have one. That makes it a
free parameter, and the audit of \S\ref{sci:S14.2} has been amended to list it as one.

\subsection{The morphisms must be logic-preserving, and that costs nothing}\label{sci:S15.3}

Functoriality on morphisms is where a real correction is needed, and we think it is the most useful
thing this section contains.

The morphisms of $\mathbf{GSLT}$ are bisimulation-preserving maps, where the bisimulation is the
context-labelled one and is a congruence relative to a class of admissible contexts
(Remark~\ref{lts:rem:relative}). But an ecology depends on \emph{spatial} structure: concealment
depth, which computations sit inside which, how a network separates when a locus is removed. And the
spatial layer is strictly more discriminating than bisimulation \emph{in the base theory} --- this is
the thesis of \cite{ClassifierNote}, that separation says things no modal logic can say. So a
bisimulation-preserving map may scramble the spatial layer, changing sensing costs and foraging
inequalities while preserving behaviour exactly.

The obvious repair is to narrow the morphism class by hand to maps preserving the generated logic in
both its layers, and note that the narrowing is strict. That is what an earlier version of this
section did, and it is the weakest paragraph it contained: a class of morphisms defined by fiat to
make an assignment functorial is not a category one has any right to expect adjoints in.

It turns out not to be a narrowing at all.

\begin{proposition}[The logic-preserving maps are bisimulation-preserving maps of the extension]
\label{sci:prop:morphisms}
Let $\GSLT$ satisfy the conditions of Theorem~\ref{ox:thm:main}. Then a map of $\GSLT$ preserves the
generated logic --- structural layer included --- exactly when it preserves context-labelled
bisimulation in the observer extension $\GSLT^{+}$. Hence the assignment of ecologies, while not
functorial on $\mathbf{GSLT}$-morphisms, is functorial on morphisms in the image of
$\Obs(-,\Sigma)$.
\end{proposition}

\begin{proof}
By Corollary~\ref{ox:cor:adequacy}, logical equivalence in the structural fragment, the equational
theory, and bisimilarity in $\GSLT^{+}$ coincide on closed terms. A map preserves an equivalence
exactly when it preserves any relation coinciding with it, so the two preservation conditions are the
same condition. Functoriality then holds on the image of $\Obs(-,\Sigma)$ by
Definition~\ref{ob:def:morphism} applied in the extended theory.
\end{proof}

This is worth more than the repair it replaces. The requirement that a morphism respect what the senses
see is not an extra demand imposed on $\mathbf{GSLT}$ from the ecology's side; it is the ordinary
requirement, made in a theory where the senses are instruments the observer holds. What looked like a
narrowing was a change of the theory the morphisms live over. And the practical consequence is that the
free and forgetful adjunctions of the machinery part are not obviously lost, since one is still in a
category of GSLTs and their bisimulation-preserving maps --- a different one.

\begin{conjecture}[Laxity]
\label{sci:conj:lax}
Even over logic-preserving morphisms the induced map on distributions is only lax. An encoding may add
rewrite steps and so inflate costs, so the pushforward distorts; the natural target is distributions
ordered by stochastic dominance rather than distributions on the nose.
\end{conjecture}

We state this and do not chase it. The three-line summary for the reader in a hurry is: a fibration, a
section for each choice of monotone cost-to-rate map and no canonical one among them, and a morphism
class that looks like a restriction of $\mathbf{GSLT}$ and is in fact $\mathbf{GSLT}$ over an
enlarged theory.

\section{Proximate and ultimate scores}\label{sci:S16}
\label{sci:sec:scores}

The original motivation for this construction measured the delta between hypothesis revisions and
punished degradation. The game supersedes that mechanism as the \emph{ultimate} criterion --- survival
is the score, and it needs no adjudicator, no referee holding a mint, and no exogenous reward signal.
Two things fall out that we would otherwise have had to impose.

\begin{itemize}
\item \textbf{Vacuity is punished by hunger, not by rule.} The hypothesis $\top$ predicts nothing,
locates no metabolic channel, yields nothing, and starves. No informativeness term is needed in any
score function, because there is no score function.
\item \textbf{Occam's razor is metabolic.} A hypothesis is worth holding exactly when expected yield
exceeds $\kappa_{\mathrm{sense}} + \kappa_{\mathrm{assay}} + \kappa_{\mathrm{break}}$ plus the cost of
holding the formula itself. Description length and food are denominated in the same currency, so
minimum description length is derived rather than stipulated.
\end{itemize}

The revision delta nevertheless survives, demoted and better placed. A scientist cannot wait until it
dies to learn which way to walk the relaxation lattice, so it needs a local signal, and
$d(\varphi_n, \varphi_{n+1})$ together with the confirm/refute counts of \S\ref{sci:sec:assay} is exactly
one. This is the proximate/ultimate distinction of behavioural ecology: viability is the ultimate
criterion, the revision delta is the proximate mechanism that tracks it, and neither is redundant.

\begin{remark}[The pragmatist's objection, and surplus]
\label{sci:rem:pragmatist}
It should be conceded plainly that this scientist is not a disinterested seeker of truth. Its
epistemology is metabolic, and truths that do not pay go unfunded. That is the standard objection to
pragmatism and we do not think it can be argued away. The construction does, however, have an unusually
concrete answer: \emph{surplus}. A scientist holding a stack larger than its foraging horizon requires
can afford inquiry that does not pay, and what it buys with the excess is curiosity --- though by
Corollary~\ref{sci:cor:kselection} it might have bought offspring instead, and the choice between them is
the life-history problem. Basic research is
a metabolic luxury, available exactly to the well-fed, and in a world of conserved tokens it is
therefore rare and worth protecting. We note this and leave the moral to others.
\end{remark}

\section{Conclusion: an embodied architecture, and the effectiveness of reason}\label{sci:S17}
\label{sci:sec:conclusion}

We close with a speculation that is not a design. Suppose a perceptual front end --- a transformer
network, wired to the world, doing what such networks are unreasonably good at --- were used to encode
sensory data as a population of computations and token stacks of the kind this note describes. We are
not proposing an interface, and nothing below depends on how the encoding is done. We are asking what
would follow if one existed, because we think the consequences bear on a question older than any of
the machinery here.

\subsection{A division of labour}\label{sci:S17.1}

The scientist of this note lives among other computations. The world does not present itself that way,
or at least does not present itself that way to us, so something must do the type conversion. That is
the front end's whole job in the proposed arrangement: not to learn, but to render the environment as
an environment of the right kind --- loci with surfaces, stacks, and channels an interaction can
address.

The division is worth stating sharply, because it inverts the usual allocation of labour:

\begin{quote}
\itshape The network proposes. The ecology is the epistemology.
\end{quote}

Nothing in \S\S\ref{sci:sec:logic}--\ref{sci:sec:repro} is a learning rule in the gradient sense, and nothing
in a perceptual encoder is a hypothesis. Keeping them apart is what makes the arrangement legible.

\begin{remark}[A correction, made in place]
\label{sci:rem:transducer}
Earlier printings of this argument gave the slogan as \emph{the network is the transducer}, and that is
wrong in a way worth recording rather than quietly repairing. If the front end is modelled as adjoining
builtin processes to the calculus, the transduction map is quotation and dereference themselves, extended
to the builtin sort: it is exact and lossless by construction, and a network can neither improve on it nor
is needed for it. What a network is needed for is the prior and much harder job of \emph{proposing} which
distinctions in the host world deserve builtin names at all. Chapter~\ref{ch:notation} states that job
precisely, gives it a success criterion --- coverage rather than reward --- and identifies causal states as
the principled first candidate. The second half of the slogan is untouched.
\end{remark}

\subsection{Why continual learning falls to the ecology}\label{sci:S17.2}

The features that make an architecture of this shape suited to continual learning are not additions to
it. They are the results already proved.

\begin{description}
\item[Revision is local.] A revision move is an operation at a locus in the formula tree
(Definition~\ref{sci:def:move}), not a global adjustment of every parameter at once. Whatever is not at the
locus is untouched, so the failure mode in which new learning overwrites old has no mechanism here.
The corresponding cost is that a locus must be found, which is what \S\ref{sci:sec:search} is about.

\item[Exploration and exploitation are a budget, not a hyperparameter.] The allocation between
foraging, inquiry, and reproduction (\S\ref{sci:sec:energy}) is a division of one conserved stack, and the
trade-off is settled by what the ecology pays for rather than by a coefficient chosen in advance.

\item[A population is required, not preferred.] This is the sharpest of the three.
Proposition~\ref{sci:prop:confinement} says an individual's incremental revision cannot leave the ball it
starts in --- not slowly, but not at all. By Remark~\ref{sci:rem:dial} the escape is recombination, which
requires more than one learner and a lineage relating them. So an architecture built on this analysis
does not have a population because populations are convenient; it has one because a single learner is
provably confined.

\item[Retirement has a criterion.] Mortality is not a scheduling policy imposed from outside. A model
that cannot fund its next step stops, and the ledger says when.
\end{description}

\subsection{The unreasonable effectiveness of reason}\label{sci:S17.3}

Wigner's question is why mathematics, developed for reasons internal to itself, turns out to describe a
world that did not consult it \cite{Wigner}. Hamming's answer is largely a selection effect: we notice
the mathematics that fits and forget the rest \cite{Hamming}. The evolutionary answer is that our
faculties were shaped by the world they model. Neither is comfortable, the first because the successes
are too systematic and too predictive, the second because the mathematics that works best --- for
spectra, for spacetime, for objects in no ancestral environment --- is exactly the mathematics no
ancestor was selected on.

The architecture posited here suggests the question has been asked across a gap that is not there.

Wigner's framing has mathematics on one side as a free creation, the world on the other as a brute
given, and an unexplained correspondence between them. In this setting there is no such separation to
explain. The hypothesis language is not chosen and then found to fit: it is generated by OSLF from the
presentation of the computational substrate itself, and adequacy is the theorem that it separates
exactly what interaction separates --- no coarser, so nothing observable is inexpressible, and no
finer, so nothing expressible is unobservable. The mathematics available to such a learner is the
mathematics of what it can distinguish, because it is the language the distinguishing relation
generates.

That sentence has one load-bearing word in it, and until recently this chapter was not entitled to it.
``Adequacy is the theorem'' presupposes that the generated logic and the observational equivalence are
calibrated against each other, and \S\ref{sci:S15.3} used to say in so many words that they are not:
that the structural layer strictly refines context-labelled bisimulation. An argument cannot rest on an
equation a later section denies. Chapter~\ref{ch:obsext} removes the denial rather than the argument
--- the two claims were about two different relations, and naming the index reconciles them
(Remark~\ref{sci:rem:adequacy-index}) --- so what follows now stands on something. It stands on less
than it appears to, and the next remark says how much less.

\begin{quote}
\itshape Mathematics is effective on exactly the domain to which interaction reaches, because it is
the shadow that the access relation casts.
\end{quote}

\begin{remark}[What this argument is standing on]
\label{sci:rem:wigner-scope}
The dissolution above is the boldest claim in this note and the reader is entitled to see the wire it hangs
from. There is no gap between the hypothesis language and the world because \S\ref{sec:isolation} removed
one: computations are ontologically isolated, so the environment of a computation is other computations,
so the logic OSLF generates over the substrate is a logic over exactly the things on the far side of the
cut. The correspondence between symbol and referent is therefore \emph{analytic} --- the ontology and the
hypothesis language are two presentations of one object --- and the effectiveness of the agent's
mathematics on its environment is a consequence of that coincidence rather than a discovery about the
environment. Within the scope, the argument is sound and unqualified. Outside it --- for an agent whose
environment is a room, a market, a coastline, another agent's body --- the coincidence lapses and the
argument requires re-derivation. Chapter~\ref{ch:trick} takes that up directly, names the move, and says
what it bought; Chapter~\ref{ch:notation} says what re-derivation would cost.
\end{remark}

Two things follow which we think are more than restatement.

First, the transfer problem dissolves in the right direction. Adequacy is not adequacy for a particular
environment; it is adequacy for interaction \emph{as such}, since context-labelled bisimulation is the
finest equivalence any interacting learner could resolve, whatever it is interacting with. That is why
the mathematics carries to domains no ancestor met. It was never fitted to the ancestral environment
in the first place --- it was fitted to the form of access, and the form of access did not change when
the subject matter did.

Second, there is a boundary, and its location is predicted rather than assumed. Effectiveness should
degrade precisely where the relation between learner and phenomenon is not one of interaction across a
cut --- where there is no context to build, no probe to install, no reply to condition on. That the
domains which have historically resisted mathematisation have tended to be of just this kind is not
proof, but it is the right shape of evidence, and it is at least a claim that could fail.

Finally, the categories themselves are accounted for rather than presupposed.
Corollary~\ref{sci:cor:induction} forces hypotheses to be about namespaces, because destructive assay
leaves an individual-level hypothesis worthless after its own first test, and
Proposition~\ref{sci:prop:kinds} sustains those namespaces through reproduction. So the traffic in kinds,
classes and universals that mathematics and science run on is not read off the world and not stipulated
by the knower. It is what a mortal, interacting learner can afford to hold. That is a Kantian
conclusion with a ledger attached, and the ledger is what makes it a claim rather than a posture.
Chapter~\ref{ch:kant} takes the remark seriously and reads the whole of Kant's transcendental
project this way.

\subsection{What is dictated, and what is not}\label{sci:S17.4}

If the mind is a computational apparatus, we think an architecture of roughly this shape is close to
forced. Take the premises separately. If a mind is computational, it is ontologically isolated, so its
environment is other computations and its access is interaction. If its access is interaction, its
hypothesis language must be adequate for context-labelled bisimulation, on pain of being blind or
wasteful. If hypotheses are to be tested rather than merely entertained, they must occupy guard
position, so that a confirmation is a rendezvous. If the apparatus is finite, its budget bounds the
fragment it can hold, and mortality supplies the motive without an external scorekeeper --- of which
none is available anyway, since a scorekeeper would be an oracle across the cut. And if revision is
search on the resulting lattice, Proposition~\ref{sci:prop:confinement} requires a population.

Each of those is a derivation in this note rather than a design decision, which is what we mean by
saying the shape is all but dictated. It should be equally clear what has not been shown. We have not
shown that the mind is computational; that premise is doing all the work and we have assumed it
throughout. We have not shown the world is computational either --- only that a computational knower's
mathematics will fit whatever it can interact with, which is a claim about the knower and not about
what is known. And we have not built the encoder, which is exactly the part where the difficulty of
any actual system would be concentrated --- a sentence that reads as a modest caveat and is not one. The
encoder is where an ontology and a grounded hypothesis language, both of which this note has been
spending freely, would have to be paid for. Chapter~\ref{ch:trick} audits that account.

What the note does claim is narrower and, we hope, more durable: that once learning is modelled as a
metered dialogue between isolated computations, the apparatus of the scientific method stops looking
like something a mind does and starts looking like something a mind, so situated, could hardly avoid.

\section{Open problems}\label{sci:S18}
\label{sci:sec:open}

\begin{enumerate}
\item \label{sci:op:noble} \textbf{Inheritance of acquired characteristics.} Remark~\ref{sci:rem:noble}
declines to read the inertness of a quotation as a germ--soma barrier, and Noble's case against a
strict barrier \cite{Noble} is a reason for caution rather than a settled verdict. The formal question
is sharp: a scientist's hypotheses are held as data and are therefore quotable and transmissible,
whereas its metabolic and sensory machinery is code. Under what pricing does a lineage transmit the
first, and what does the resulting two-channel inheritance do to
Proposition~\ref{sci:prop:sexdefence}?

\item \textbf{The plagiarism equilibrium.} Reading a rival's model is cheaper than experimenting and
kills nobody, but holding a model requires exposing it as a guard. Under what pricing does honest
science survive theory theft, and is there a mixed equilibrium?

\item \textbf{Co-learning and non-stationarity.} When the environment is mostly scientists, the target
of every hypothesis is itself revising. Convergence becomes a fixed-point question rather than a
learning-theoretic one, and this is where we expect the geometry/matter corecursion to be readable as
learning.

\item \textbf{Laxity.} Conjecture~\ref{sci:conj:lax}: is the induced map on distributions lax, and is
stochastic dominance the right order on the target?

\item \textbf{The revision policy.} We left the agent design open on purpose. Is there a policy on the
relaxation lattice that is optimal in the ecological sense --- maximising expected residual life
rather than predictive accuracy --- and does it look anything like the policies learning theory
recommends?

\item \textbf{Is the mixed diet dominated?} Propositions~\ref{sci:prop:worstofboth}
and~\ref{sci:prop:composite} are the two halves of a claim about trophic sparsity. Testing them in the
ensemble of \S\ref{sci:sec:distributions} --- is the mixed access profile ever optimal for an atom, and
does the composite dominate it where it is not --- is the first computation we would run.

\item \textbf{Sharpening the phase boundary.} Figure~\ref{sci:fig:phase} is drawn from the sign of a cost
inequality. Making its boundaries quantitative for a specific pricing schedule, via the sampling
construction of Chapter~\ref{ch:weight}, is the first computation to do.
\end{enumerate}


\section{Appendix: A worked ecosystem}\label{sci:S19}
\label{sci:app:eco}

We give a small world containing two large external token stacks, three grazing non-scientist mortal
computations, two reproductively isolated species of sexually reproducing scientist, and one asexual
lineage. It is written in core rholang extended with the \texttt{where} clause of \S\ref{sci:sec:assay};
stacks are integers and arithmetic is taken as given. The point of the listing is not that it is
efficient but that every construction of the note appears in it as ordinary code, with no apparatus
added.

\subsection{Sources, metabolism, death}\label{sci:S19.1}

\begin{lstlisting}
// A source (Def. "Source and prey"): rate-limited by being a single contract, non-exclusive
// since anyone may call it, non-adaptive, and FINITE -- conservation is exact.
def Source(s, q, reserve) = {
  for (ret <- s) {
    match reserve >= q {
      true  => { ret!(q) | Source(s, q, reserve - q) }
      false => { ret!(0) }                  // exhausted; cf. Prop. "Trophic succession"
    }
  }
}

// Reclaim - debit - reissue (Def. "Metabolic channel").  Between the receipt and the send the
// stack is a message in flight: metabolism is a race condition (Rem. "Metabolism is a race condition"), and
// the metabolic channel m is exactly the vulnerability of Thm. "Predators are non-image contexts".
def Live(m, burn, step) = {
  for (sigma <- m) {
    match *sigma > burn {
      true  => { m!(*sigma - burn) | *step | Live(m, burn, step) }
      false => { Nil }                      // death is failure to fund a step
    }
  }
}
\end{lstlisting}

\subsection{A non-scientist}\label{sci:S19.2}

\begin{lstlisting}
// A grazer holds no hypothesis and senses nothing beyond the one channel it
// was born knowing.  By Prop. "The two sciences" its environment is a source, so its science
// was over before it began -- and it pays no epistemic rent at all.
def Grazer(m, src) = {
  Live(m, 1, @{
    new probe in {
        src!(*probe)
      | for (g <- probe) { for (sigma <- m) { m!(*sigma + *g) } }
    }
  })
}
\end{lstlisting}

\subsection{Foraging as an assay}\label{sci:S19.3}

\begin{lstlisting}
// The complementary guard pair of Def. "Assay", verbatim.  The hypothesis is not
// consulted and then acted on; it sits in guard position, so confirmation IS
// the rendezvous.  If the source is exhausted neither guard fires: that is
// the third outcome of Prop. "Trichotomy", and it is not the same as refutation.
def Forage(m, sns, hyp) = {
  new probe in {
      sns!(*probe)                                 // the probe perturbs
    | for (g <- probe where g |= *hyp) {           // predicted
        for (sigma <- m) { m!(*sigma + *g) }
      }
    | for (g <- probe where ~(g |= *hyp)) {        // refuted
        Revise(hyp)                                // walk the lattice (Sec. 5)
      }
  }
}
\end{lstlisting}

\subsection{Birth: a quotation, then a drop}\label{sci:S19.4}

\begin{lstlisting}
// The offspring's code is CONSTRUCTED as a name and sits inert on `code'
// -- it does not reduce, does not age, burns nothing (Def. "Genome, phenotype").  Dropping it
// with * is what makes it a phenotype, and what starts the meter.
//
// The guard enforces Prop. "Minimum viable endowment": a parent may not endow below the outlay its
// child needs for a first successful harvest.
def Beget(m, alleles, prv, mate) = {
  for (sigma <- m where *sigma > *prv + floor) {
    m!(*sigma - *prv)                             // the endowment is subtracted
  | match *alleles {
      [s, h, p] => {
        new childM, code in {
            code!( @{ Body(childM, s, h, p, mate) } )      // the germ: inert
          | for (g <- code) { *g | childM!(*prv) }         // the soma: metered
        }
      }
    }
  }
}
\end{lstlisting}

\subsection{Crossover}\label{sci:S19.5}

\begin{lstlisting}
// Three homologous loci are exposed as parameters: the sensing allele, the
// hypothesis allele, and the provisioning allele.  This is assortment at
// three loci; crossover proper (Sec. "Crossover") splits a guard from its body
// inside an allele and carries the reading-frame condition of Def. "Homology".
def Pick(a, b, bit, ret) = {
  match *bit { 0 => { ret!(*a) } ; 1 => { ret!(*b) } }
}

def Crossover(mother, father, chi, ret) = {
  match (*mother, *father, *chi) {
    ([sm,hm,pm], [sf,hf,pf], [b1,b2,b3]) => {
      new r1, r2, r3 in {
          Pick(sm,sf,b1,r1) | Pick(hm,hf,b2,r2) | Pick(pm,pf,b3,r3)
        | for (s <- r1 & h <- r2 & p <- r3) { ret!([*s, *h, *p]) }
      }
    }
  }
}
\end{lstlisting}

\subsection{A sexual scientist, and an asexual one}\label{sci:S19.6}

\begin{lstlisting}
// Mating requires DISCLOSURE: the genome goes out on a public channel, where
// anything able to read it can.  Conspecific(s) is a name predicate (Sec. 3.1)
// holding of genomes whose sensing allele names the same source namespace as s.
// Nothing declares a species; reproductive isolation is namespace disjointness
// and nothing more (Prop. "Species is a namespace formula").
def Mate(m, alleles, sns, prv, mate) = {
  mate!(*alleles)
| for (partner <- mate where partner |= Conspecific(*sns)) {
    new chi, ret in {
        chi!( Coin3 )                        // the choice function of Def. "Crossover"
      | Crossover(alleles, partner, chi, ret)
      | for (kid <- ret) { Beget(m, kid, prv, mate) }
    }
  }
}

// The shared developmental program.  Alleles vary; this does not.
def Body(m, sns, hyp, prv, mate) = {
  Live(m, 2, @{                              // scientists burn faster than grazers
      Forage(m, sns, hyp)
    | Mate(m, @[*sns, *hyp, *prv], sns, prv, mate)
  })
}

// The asexual lineage: no partner, no alignment, NO DISCLOSURE, and no
// recombination to compute.  Strictly cheaper -- and by Prop. "Sex attacks the predator's hypothesis language" it presents
// a fixed namespace for anything that cares to learn it.
def Clone(m, sns, hyp, prv) = {
  Live(m, 2, @{
      Forage(m, sns, hyp)
    | Beget(m, @[*sns, *hyp, *prv], prv, @Nil)
  })
}
\end{lstlisting}

\subsection{The world}\label{sci:S19.7}

\begin{lstlisting}
new sun1, sun2, mate in {

  // Two large external stacks.  sun1 is shallow and gives small quanta;
  // sun2 lies deeper and gives larger ones -- the two order parameters
  // of Sec. 13.2 instantiated as two loci rather than as a distribution.
    Source(sun1, 1,  60000)
  | Source(sun2, 8, 240000)

  // Three non-scientists.
  | new g1, g2, g3 in {
        g1!(50) | Grazer(g1, sun1)
      | g2!(50) | Grazer(g2, sun1)
      | g3!(80) | Grazer(g3, sun2)
    }

  // Species A: hunts sun1, cheap hypothesis, many small offspring.
  | new a1, a2 in {
        a1!(400) | Body(a1, @sun1, @{ q => q > 0 }, @20, mate)
      | a2!(400) | Body(a2, @sun1, @{ q => q > 0 }, @20, mate)
    }

  // Species B: hunts sun2, a discriminating hypothesis, few large offspring.
  // A and B share the mating channel and never interbreed.
  | new b1, b2 in {
        b1!(900) | Body(b1, @sun2, @{ q => q > 4 }, @150, mate)
      | b2!(900) | Body(b2, @sun2, @{ q => q > 4 }, @150, mate)
    }

  // The asexual lineage, competing with species A for sun1.
  | new c1 in { c1!(500) | Clone(c1, @sun1, @{ q => q > 0 }, @60) }
}
\end{lstlisting}

\subsection{What to watch}\label{sci:S19.8}

\begin{itemize}
\item \textbf{Conservation is checkable by inspection.} $\Theta$ is the sum of the two reserves and
the eight initial stacks. Nothing in the listing mints; \texttt{Source} only ever moves reserve into a
consumer, and \texttt{Beget} only ever moves a parent's stack into a child's. Every debit taken by
\texttt{Live} is a token leaving $\sigma$ for $\kappa$, so $\sigma + \kappa = \Theta$ throughout.

\item \textbf{The three outcomes are structural, not bookkeeping.} In \texttt{Forage}, refutation is
an event because $\lnot\varphi$ is a guard; the case where the source is exhausted produces neither
event. Proposition~\ref{sci:prop:trichotomy} is visible as three code paths, one of which is empty.

\item \textbf{Nothing screens.} No embryo is model-checked anywhere in the listing. A parent could
in principle check its recombinant against the viability formula before dropping it, and does not,
because the check is priced against the same stack --- Proposition~\ref{sci:prop:selection} as a
programming decision rather than an argument.

\item \textbf{Speciation costs one guard.} $A$ and $B$ share \texttt{mate} and are nevertheless
isolated. Delete the \texttt{where} clause from \texttt{Mate} and they hybridise; the offspring
inherits a sensing allele naming one source and a hypothesis allele tuned to the other, and starves
without ever being refuted. That is the reading-frame condition failing at the level of the ecology.

\item \textbf{The race is real.} \texttt{Live} and \texttt{Forage} both take from \texttt{m} and put
back. They conserve the stack and cannot deadlock, but they interleave --- which is
Remark~\ref{sci:rem:race}'s point that living means racing others to your own stack, with the others
here being your own body.

\item \textbf{This world is at the autotrophic end.} Every organism feeds on a source, so by
Proposition~\ref{sci:prop:twosciences} every lineage's inquiry terminates and none pays an unbounded
epistemic rent. Note how close the other end is: species $B$'s sensing allele already names a deeper
namespace, and a single reallocation of it toward a grazer's \texttt{m} converts the lineage to
heterotrophy, with Corollary~\ref{sci:cor:intelligence}'s standing cost arriving in the same step.

\item \textbf{And it ends.} When both reserves reach zero, \texttt{Source} returns $0$ forever, no
guard on a positive quantum fires again, and the remaining tokens sit in metabolic channels. From
there the only credit available is another organism's stack. Proposition~\ref{sci:prop:succession} is not
an asymptotic remark about this listing; it is the last thing that happens in it.
\end{itemize}

